% ============================================================
%  第二卷 第十四章 相对论宇宙学（§111–§119 + 译后记）
%  底本：高教社中文版第8版；OCR 错误已修正，脚注文本待脚注代理补录
%  主要修正：
%  - §111："其中4是一个新的常数"→Λ；附加项 Ag_ik→Λg_ik；
%    "如果用“角度∵χ”"→“角度”χ；"在这些坐标中的球的表面积是 χ aχ 4πa²sin²χ"
%    按 p.401 剔除 OCR 乱码，仅保留 4πa²sin²χ；
%  - §112：克里斯托费尔符号 Γ^i_{kl} OCR 误为 T^i_{kl}，恢复；"三维张量一 P"
%    之"一"为负号，复原为 −P_α^β；
%  - §113：弗里德曼型 exercize 引文 OCR 残缺"(E. M., . . , . . , 19)"按 p.409
%    复原为（Е. М. Лифшиц, И. М. Халатников, 1960）；场方程 8πk/c⁴→8πk/c²
%    （(113.10) 后一式，按 p.409）；
%  - §114：(114.18) "I = cos t"→const；习题2 积分上限 x→χ，句尾衍字"5"删；
%  - §115：(115.5) tag OCR 残缺复原；(115.19)—(115.22) 条件短语
%    "⊥/stackrel{\perp}{=}、X f ////、\Y H"等乱码按 p.421–423 复原为"当/对于条件
%    p=ε/3、当 p=0"；"按照 a⁻² 减AD②"→"减小②"；
%    §113 习题解中 "√n~1/t"、"√n 它的数值"按 p.410 复原为"而 ∼1/t 的项"、
%    "而它的数值"；"度规 α_αβ"→a_{αβ}；
%  - §116：比安基分类表 OCR 全乱（WI_0/M_0/II(a=1)W_a 等），按 p.428 表格复原
%    （I, II, VII_0, VI_0, IX, VIII, V, IV, VII_a, III(a=1)/VI_a(a≠1)）；
%    (116.8a) 一段与 ξ 次段系 p.425 脚注①正文，按约定移除，标记保留；
%    (116.24) 第三项乱排的 "=" 复原为 "−"；(116.26) 末式导数点位置按
%    (97.13) 体例复原为 (√η κ)^˙ − P（P 不求导）；
%  - §117：习题2 末度规 ±2x^{2p₂}dx²dx³ 前的 "±" 按原书保留；
%  - §118：标题"靠近奇点的振动状态"；"以第类"→"以第IX类"；
%    "结构常数如 F①"→"结构常数如下①："；标架矢量 l,m,n 显式式及体积式
%    系 p.434 脚注①正文，按约定移除，标记保留；引文按 p.434 复原为
%    （В. А. Белинский, Е. М. Лифшиц, И. М. Халатников, 1968）；"演化#①"→"演化①"；
%  - §119："物理奇点物质密度"补破折号"——"；
%  - 脚注圈码①②③照 OCR 保留；星号脚注（§111、§114 各一处）以 \textsuperscript{*}
%    保留标记；OCR 中外溢的脚注正文块（§113 两处、§116 三处、§118 两处等）
%    按约定移除，由脚注代理统一补录；
%  - 插图：图25（§118，p450）。
%  - 原书正体关键词（各向同性宇宙模型、封闭模型、开放模型、红移、哈勃常数、
%    运动群、结构常数、雅可比恒等式等）为黑体强调，OCR 无字体信息，未加粗。
%  - 2026-10-10 脚注已转录：33 处圈码标记全部替换为 \footnote{...}；3 处星号脚注
%    以 \textsuperscript{*} 标记 + 局部 * 编号的 \footnotetext 补录（§111 p.399 一处、
%    §114 p.413/p.415 两处）。原书 p.399 页脚为 ①/*/② 三条，OCR 所缺之第三条即
%    星号中译注，并非 ③；§114 星号脚注实为两处（原书如此）。
% ============================================================
\chapter{相对论宇宙学}

\section{各向同性空间}

广义相对论开辟了在宇宙学尺度上研究并解决宇宙性质问题的新途径.由此产生的新奇的可能现象是与时空的非伽利略性质相联系的（由爱因斯坦首先在1917年指出）.

这些可能现象更本质的意义在于，牛顿力学在这里会遇到矛盾的结果，这些矛盾在形式足够普遍的非相对论理论范围内是不能绕过的.譬如说，平直的（在牛顿力学中就是这样）无穷大空间被任意分布、在任何地方都不会消失、具有一定平均密度的物质所填充，那么当用牛顿力学公式计算其中的引力势时，我们会发现在每一点的引力势都趋向无穷大.这会导致作用在物质上的力为无穷大，就是说，导致悖论.

在着手系统地建立相对论宇宙学模型之前，我们对作为出发点的基本场方程作下列说明.

在§93中给出了作为确定引力场的作用量所需满足的条件，在给标量 $G$ 加上一个常数项之后，该条件仍旧能得到满足，就是说，令
\begin{equation*}
  S_{\mathrm{g}}=-\frac{c^{3}}{16\pi k}\int(G+2\Lambda)\sqrt{-g}\,\mathrm{d}\varOmega,
\end{equation*}
其中 $\Lambda$ 是一个新的常数（带有量纲 $\mathrm{cm}^{-2}$）.这样的改变会导致爱因斯坦方程中出现一个附加项 $\Lambda g_{ik}$：
\begin{equation*}
  R_{ik}-\frac{1}{2}Rg_{ik}=\frac{8\pi k}{c^{4}}T_{ik}+\Lambda g_{ik}.
\end{equation*}

如果赋予“宇宙学常数”$\Lambda$ 一个很小的值，那么这个项的出现对于不太大的时空区域内的引力场将不会造成显著影响，但是会导致新类型的、或许能够在整体上描述宇宙的“宇宙学解”的出现\footnote{其中会出现稳态解，而当 $\Lambda=0$ 时则不存在这些稳态解. 正是出于这个目的，爱因斯坦引入了“宇宙学项”. 这个情况发生在弗里德曼发现场方程的非稳态解之前——见下文.}.然而，在当今时代，对于基本理论方程在形式上的这种改变，无论在观测方面，还是在理论方面，都没有任何坚实的和令人信服的根据.我们强调，这里所说的是具有深刻物理涵义的改变：向拉格朗日函数的密度中引入根本不依赖于场的状态的常数项，这意味着给时空赋予一个原则上不可消除的曲率，这个曲率既与物质无关，又与引力波无关.因此，本章中所有以下的叙述都是基于“经典”形式的爱因斯坦方程，而不考虑宇宙学常数.\textsuperscript{*}{\renewcommand{\thefootnote}{*}\footnotetext{高红移超新星的观测显示宇宙正在加速膨胀，为 $\Lambda$ 不等于零或性质未知的暗能量存在的可能性提供了新的证据，参见 S. Perlmutter, et al., 1997.——中译注}}

众所周知，恒星以非常不均匀的形式分布于空间——它们集中在分立的恒星系统（星系）中.但是在“大尺度”上研究宇宙时，应该忽略物质在恒星和星系中聚集而引起的“局部”非均匀性.因此，质量密度应该理解为：在线度大于星系之间距离的空间区域内的平均密度.

以下（§111—§114）研究的爱因斯坦方程的解被称作各向同性宇宙模型（由 A. A. 弗里德曼在1922年首先发现），是基于物质沿空间分布的均匀性和各向同性的假定.现有的天文学数据与这种假定并不矛盾\footnote{这里指的是关于星系在空间中的分布数据和背景射电辐射各向同性的数据.}，并且从现有的一切证据可以认为，各向同性模型大体上不但对现在的宇宙能够给出适当的描述，而且对宇宙过去的演化过程中的相当一部分也是如此.我们在下文将看到，这个模型的基本性质是它的非稳态性.毋庸置疑，这个性质（“膨胀的宇宙”）能够对宇宙学的基础问题——红移现象给出正确的解释（§114）.

同时可以明白，关于宇宙的均匀性和各向同性的假定就其自身的本质而言，不可避免地只能具有近似的特点，因为在过渡到更小的尺度时，这些性质必然会被破坏.关于宇宙的非均匀性在宇宙学问题的各个方面中可能起到的作用这个课题，我们将在§115—§119节中讨论.

空间的均匀性和各向同性意味着能够选择这样的世界时间，使得在它的每一时刻，空间的度规在所有的点上和在所有的方向上都是一样的.

首先我们来研究各向同性空间的度规本身，暂时不考虑它与时间可能有的依赖关系.如同我们在前文中的做法，将三维度规张量表示为 $\gamma_{\alpha\beta}$，就是说，将空间距离元写成如下形式：
\begin{equation}
  \mathrm{d}l^{2}=\gamma_{\alpha\beta}\mathrm{d}x^{\alpha}\mathrm{d}x^{\beta}.
\end{equation}

空间的曲率完全由空间的三维曲率张量所决定，我们将它记作 $P_{\alpha\beta\gamma\delta}$ 以区别于四维张量 $R_{iklm}$.在完全各向同性的情况下，张量 $P_{\alpha\beta\gamma\delta}$ 显然应当只用度规张量 $\gamma_{\alpha\beta}$ 来表示.因此，从自身的对称性很容易看出，它应当有下面的形式
\begin{equation}
  P_{\alpha\beta\gamma\delta}
  =\lambda(\gamma_{\alpha\gamma}\gamma_{\beta\delta}
  -\gamma_{\alpha\delta}\gamma_{\beta\gamma}),
\end{equation}
其中，$\lambda$ 是某一常数.相应地，里奇张量 $P_{\alpha\beta}=P_{\alpha\gamma\beta}^{\gamma}$ 等于
\begin{equation}
  P_{\alpha\beta}=2\lambda\gamma_{\alpha\beta},
\end{equation}
而曲率标量
\begin{equation}
  P=6\lambda.
\end{equation}

这样一来，各向同性空间的曲率特性仅用一个常数来决定.与此相应，对于空间度规总共可能有三个重要的不同情形：（1）所谓恒定正曲率空间（与正的 $\lambda$ 值相应），（2）恒定负曲率空间（与 $\lambda<0$ 的情形相应），（3）零曲率空间（与 $\lambda=0$ 的情形相应）.最后一个当然是平直空间，即欧氏空间.

为了研究度规，最便利的是从几何的相似出发，将各向同性的三维空间的几何看做是一个在假想的四维空间内、已知其为各向同性的超曲面上的几何\footnote{这个四维空间应该理解为与四维时空没有关系.}.这样的曲面是一个超球；与它相应的三维空间是恒定正曲率空间.四维空间 $x_{1},x_{2},x_{3},x_{4}$ 内半径为 $a$ 的超球的方程如下：
\begin{equation*}
  x_{1}^{2}+x_{2}^{2}+x_{3}^{2}+x_{4}^{2}=a^{2},
\end{equation*}
在其上的线元可以表示为
\begin{equation*}
  \mathrm{d}l^{2}=\mathrm{d}x_{1}^{2}+\mathrm{d}x_{2}^{2}
  +\mathrm{d}x_{3}^{2}+\mathrm{d}x_{4}^{2}.
\end{equation*}

将 $x_{1},x_{2},x_{3}$ 看做是三个空间坐标，利用第一个方程，从 $\mathrm{d}l^{2}$ 中消去假想坐标 $x_{4}$，我们得到空间距离元如下：
\begin{equation}
  \mathrm{d}l^{2}=\mathrm{d}x_{1}^{2}+\mathrm{d}x_{2}^{2}+\mathrm{d}x_{3}^{2}
  +\frac{(x_{1}\mathrm{d}x_{1}+x_{2}\mathrm{d}x_{2}+x_{3}\mathrm{d}x_{3})^{2}}
  {a^{2}-x_{1}^{2}-x_{2}^{2}-x_{3}^{2}}.
\end{equation}

从这个式子不难计算（111.2）中的常数 $\lambda$.既然我们早已知道 $P_{\alpha\beta}$ 在整个空间都有（111.3）的形式，那么，只须计算它在原点附近的一点上的值就够了，在这一点上，$\gamma_{\alpha\beta}$ 等于
\begin{equation*}
  \gamma_{\alpha\beta}=\delta_{\alpha\beta}+\frac{x_{\alpha}x_{\beta}}{a^{2}}.
\end{equation*}

因为 $\gamma_{\alpha\beta}$ 的一阶导数，从而量 $\lambda_{\beta\gamma}^{\alpha}$（比较§88节，习题1）——对应于度规 $\gamma_{\alpha\beta}$ 的三维克里斯托夫符号——在坐标原点为零，所以根据通式（92.7）来计算是很简单的，结果得到
\begin{equation}
  \lambda=\frac{1}{a^{2}}.
\end{equation}

我们可以称 $a$ 为空间的曲率半径.引入相应的“球”坐标 $r,\theta,\varphi$ 来代替坐标 $x_{1},x_{2},x_{3}$.这时，线元的表达式将有如下的形式：
\begin{equation}
  \mathrm{d}l^{2}=\frac{\mathrm{d}r^{2}}{1-r^{2}/a^{2}}
  +r^{2}(\sin^{2}\theta\,\mathrm{d}\varphi^{2}+\mathrm{d}\theta^{2}).
\end{equation}

坐标原点当然可以选择在空间内任何一点.在这些坐标中，圆的周长是 $2\pi r$ 而球的表面积是 $4\pi r^{2}$.圆（或球）的“半径”等于
\begin{equation*}
  \int_{0}^{r}\frac{\mathrm{d}r}{\sqrt{1-r^{2}/a^{2}}}=a\arcsin\frac{r}{a},
\end{equation*}
即大于 $r$.因此，在这个空间中，圆周与半径之比将小于 $2\pi$.

如果用“角度”$\chi$ 按照 $r=a\sin\chi$（$\chi$ 的变化范围是从0到 $\pi$）的关系来代替坐标 $r$，就能得到一个四维球坐标\footnote{笛卡儿坐标 $x_{1},x_{2},x_{3},x_{4}$ 与四维球坐标 $a,\theta,\varphi,\chi$ 有下面的关系：
\[
x_{1}=a\sin\chi\sin\theta\cos\varphi,\qquad
x_{2}=a\sin\chi\sin\theta\sin\varphi,
\]
\[
x_{3}=a\sin\chi\cos\theta,\qquad
x_{4}=a\cos\chi.
\]}，于是就可以写出 $\mathrm{d}l$ 的另一个便利的形式：
\begin{equation}
  \mathrm{d}l^{2}=a^{2}[\mathrm{d}\chi^{2}
  +\sin^{2}\chi(\sin^{2}\theta\,\mathrm{d}\varphi^{2}+\mathrm{d}\theta^{2})].
\end{equation}

坐标 $\chi$ 度量到原点的距离，这个距离为 $a\chi$.在这些坐标中的球的表面积是 $4\pi a^{2}\sin^{2}\chi$.我们看出，当我们从坐标原点离开时，球的面积随之而增加，当距离原点 $\pi a/2$ 时，球的面积达到了最大值 $4\pi a^{2}$.此后，这个面积开始减小，在空间的“对立极点”，与原点相距为 $\pi a$（在这样的空间内它是一般可能存在的最大距离），球的面积化为一点（所有这些，只要我们注意到坐标 $r$ 不能取大于 $a$ 的值，就可以从（111.7）看出）.

一个有正曲率的空间的体积等于
\begin{equation*}
  V=\int_{0}^{2\pi}\int_{0}^{\pi}\int_{0}^{\pi}
  a^{3}\sin^{2}\chi\sin\theta\,\mathrm{d}\chi\,\mathrm{d}\theta\,\mathrm{d}\varphi,
\end{equation*}
由此可得
\begin{equation}
  V=2\pi^{2}a^{3}.
\end{equation}

因此，一个有正曲率的空间是“自封闭的”，它的体积是有限的，但是，不言而喻，它没有边界.

值得指出，在封闭空间中，总电荷必须是零.事实上，在一个有限空间中，每个封闭曲面在它自身的两边都包围着空间的一个有限区域.因此，一方面，电场经过这个曲面的通量等于在这个曲面内的总电荷，而另一方面，等于曲面之外的总电荷，只是正负号相反.因此，这个曲面两边的电荷之和为零.

类似地，从四维动量的曲面积分形式的表达式（96.16）可知，整个空间中的四维总动量 $P^{i}$ 为零.

现在我们来研究有负的恒定曲率的空间的几何.从（111.6）我们看到，如果 $a$ 是虚数，那么，$\lambda$ 就是负的.因此，对于负曲率空间的所有公式，只须用 $\mathrm{i}a$ 代 $a$，就立即可以从前面的公式得出.换句话说，负曲率空间的几何在数学上可看做在一个半径为虚数的四维伪球上的几何.

因此，常数 $\lambda$ 现在等于
\begin{equation}
  \lambda=-\frac{1}{a^{2}}.
\end{equation}

负曲率空间中的线元在 $r,\theta,\varphi$ 坐标中有下面的形式：
\begin{equation}
  \mathrm{d}l^{2}=\frac{\mathrm{d}r^{2}}{1+r^{2}/a^{2}}
  +r^{2}(\sin^{2}\theta\,\mathrm{d}\varphi^{2}+\mathrm{d}\theta^{2}),
\end{equation}
其中，$r$ 可以取从0到 $\infty$ 之间的所有值.圆的周长与半径之比现在大于 $2\pi$.如果按照 $r=a\sinh\chi$（$\chi$ 从0到 $\infty$）引入坐标 $\chi$，我们就得到与（111.8）相应的 $\mathrm{d}l^{2}$ 的表达式
\begin{equation}
  \mathrm{d}l^{2}=a^{2}\{\mathrm{d}\chi^{2}
  +\sinh^{2}\chi(\sin^{2}\theta\,\mathrm{d}\varphi^{2}+\mathrm{d}\theta^{2})\}.
\end{equation}

球的面积现在等于 $4\pi a^{2}\sinh^{2}\chi$，当我们从原点移开时（$\chi$ 因之增加），这个面积将无限制地增加.负曲率空间的体积显然是无限的.

\begin{xiti}

\noindent{\heiti 习题}\quad 将线元（111.7）变换成这样的形式，在这个形式中，线元与其欧几里得表达式成比例（共形欧氏坐标）.

解：将
\begin{equation*}
  r=\frac{r_{1}}{1+\dfrac{r_{1}^{2}}{4a^{2}}}
\end{equation*}
代入，得到
\begin{equation*}
  \mathrm{d}l^{2}=\left(1+\frac{r_{1}^{2}}{4a^{2}}\right)^{-2}
  (\mathrm{d}r_{1}^{2}+r_{1}^{2}\mathrm{d}\theta^{2}
  +r_{1}^{2}\sin^{2}\theta\,\mathrm{d}\varphi^{2}).
\end{equation*}

\end{xiti}

\section{封闭的各向同性模型}

为了研究各向同性模型的时空度规，首先必须选定参考系.最便利的参考系是这样的一个“共动”参考系，它在空间的每一点随着在该点的物质一起运动.换句话说，这个参考系恰恰就是充满空间的物质；根据定义，物质在这个参考系中的速度处处都为零.显而易见，参考系的这种选择对于各向同性的模型是合理的；作任何其他的选择时，物体速度的方向就造成空间的不同方向在外表上看起来不等价.时间的坐标应当像上节开始所说的那样选择，就是说，使得在每一时刻，度规在整个空间内都是一样的.

由于所有方向完全等价，度规张量的分量 $g_{0\alpha}$ 在我们所选择的参考系内等于零.事实上，如果三个分量 $g_{0\alpha}$ 不为零，则它们可以当做一个三维矢量的分量，那么，不同的方向就不等价了.因此 $\mathrm{d}s^{2}$ 应当有 $\mathrm{d}s^{2}=g_{00}(\mathrm{d}x^{0})^{2}-\mathrm{d}l^{2}$ 的形式.分量 $g_{00}$ 在这里仅仅是 $x^{0}$ 的函数.于是，我们总能够通过选择时间坐标以使得 $g_{00}$ 化为1.用 $ct$ 表示这样选择的时间坐标，我们得到
\begin{equation}
  \mathrm{d}s^{2}=c^{2}\mathrm{d}t^{2}-\mathrm{d}l^{2}.
\end{equation}
变量 $t$ 是空间中每一点的同步的固有时.

我们从研究正曲率空间开始；为了简便起见，下面我们将爱因斯坦方程的相应的解称为封闭模型.对于 $\mathrm{d}l$，我们用表达式（111.8），在其中的曲率半径 $a$ 一般来说是时间的函数.因此，我们将 $\mathrm{d}s^{2}$ 写成
\begin{equation}
  \mathrm{d}s^{2}=c^{2}\mathrm{d}t^{2}
  -a^{2}(t)\{\mathrm{d}\chi^{2}
  +\sin^{2}\chi(\mathrm{d}\theta^{2}+\sin^{2}\theta\,\mathrm{d}\varphi^{2})\}.
\end{equation}

函数 $a(t)$ 由爱因斯坦方程所决定.为了解这些方程，用由关系式
\begin{equation}
  c\,\mathrm{d}t=a\,\mathrm{d}\eta
\end{equation}
定义的 $\eta$ 来代替时间是便利的.这时，$\mathrm{d}s^{2}$ 可以写成
\begin{equation}
  \mathrm{d}s^{2}=a^{2}(\eta)\{\mathrm{d}\eta^{2}-\mathrm{d}\chi^{2}
  -\sin^{2}\chi(\mathrm{d}\theta^{2}+\sin^{2}\theta\,\mathrm{d}\varphi^{2})\}.
\end{equation}

要建立场方程，应从计算张量 $R_{ik}$ 的分量开始（$\eta,\chi,\theta,\varphi$ 是坐标 $x^{0},x^{1},x^{2},x^{3}$）.利用度规张量的分量的值
\begin{equation*}
  g_{00}=a^{2},\quad g_{11}=-a^{2},\quad
  g_{22}=-a^{2}\sin^{2}\chi,\quad
  g_{33}=-a^{2}\sin^{2}\chi\sin^{2}\theta
\end{equation*}
计算 $\Gamma_{kl}^{i}$ 诸量：
\begin{equation*}
  \Gamma_{00}^{0}=\frac{a}{a},\quad
  \Gamma_{\alpha\beta}^{0}=-\frac{a}{a^{3}}g_{\alpha\beta},\quad
  \Gamma_{0\beta}^{\alpha}=\frac{a}{a}\delta_{\beta}^{\alpha},\quad
  \Gamma_{\alpha 0}^{0}=\Gamma_{00}^{\alpha}=0,
\end{equation*}
其中撇号表示对 $\eta$ 微分（分量 $\Gamma_{\beta\gamma}^{\alpha}$ 的表达式没有必要计算出来）.利用这些值，按照通式（92.7），我们得到
\begin{equation*}
  R_{0}^{0}=\frac{3}{a^{4}}(a^{2}-aa).
\end{equation*}

出于对称性的考虑（就如上文对 $g_{0\alpha}$ 所做的），我们预先断定，分量 $R_{0\alpha}=0$.对于分量 $R_{\alpha}^{\beta}$ 的计算，我们指出，如果在它们当中分离只包含 $g_{\alpha\beta}$ 的那些项（也就是说，只有 $\Gamma_{\beta\gamma}^{\alpha}$），那么这些项应该构成三维张量 $-P_{\alpha}^{\beta}$ 的分量，这些分量的值从（111.3）和（111.6）就已经得知了：
\begin{equation*}
  R_{\alpha}^{\beta}=-P_{\alpha}^{\beta}+\cdots
  =-\frac{2}{a^{2}}\delta_{\alpha}^{\beta}+\cdots,
\end{equation*}
其中省略号指的是那些同时包含 $g_{\alpha\beta}$ 和 $g_{00}$ 的项.通过对上式的计算，我们得到：
\begin{equation*}
  R_{\alpha}^{\beta}
  =-\frac{1}{a^{4}}(2a^{2}+a^{2}+aa)\delta_{\alpha}^{\beta},
\end{equation*}
然后
\begin{equation*}
  R=R_{0}^{0}+R_{\alpha}^{\alpha}
  =-\frac{6}{a^{3}}(a+a).
\end{equation*}

既然在我们所选择的参考系中物质是静止的，那么，$u^{\alpha}=0,\ u^{0}=1/a$，于是从（94.9）式得到 $T_{0}^{0}=\varepsilon$，其中 $\varepsilon$ 是物质的能量密度.将得到的关系式代入方程
\begin{equation*}
  R_{0}^{0}-\frac{1}{2}R=\frac{8\pi k}{c^{4}}T_{0}^{0},
\end{equation*}
我们得到
\begin{equation}
  \frac{8\pi k}{c^{4}}\varepsilon=\frac{3}{a^{4}}(a^{2}+a^{2}).
\end{equation}

这里出现了两个未知函数 $\varepsilon$ 和 $a$；因此，我们必须还要找到另外一个方程.为此，选择方程 $T_{0;i}^{i}=0$ 是便利的（用来代替爱因斯坦方程的空间分量），这个方程是四个方程（94.7）之中的一个，如我们所知，它是包含在场方程之内的.这个方程也可利用热力学关系用下面的方法直接导出.

在场方程中应用能量动量张量的表达式（94.9）时，我们省略了所有导致熵增加的能量耗散过程.这个省略在这里当然是完全合理的，因为由于能量耗散而应该加到 $T_{k}^{i}$ 上的一些附加项与能量密度 $\varepsilon$ 相比是微不足道的，这里最后提到的能量密度包含了物体的静能.

因此，在推导场方程时，我们可以将总熵当做是不变的.现在我们来应用已知的热力学关系式 $\mathrm{d}\mathcal{E}=T\mathrm{d}S-p\,\mathrm{d}V$，此处的 $\mathcal{E},S,V$ 是体系的能量、熵与体积，而 $p,T$ 则是它的压强与温度.在熵不变的情况下，我们简单地有 $\mathrm{d}\mathcal{E}=-p\,\mathrm{d}V$.引入能量密度 $\varepsilon=\mathcal{E}/V$，我们很容易求出
\begin{equation*}
  \mathrm{d}\varepsilon=-(\varepsilon+p)\frac{\mathrm{d}V}{V}.
\end{equation*}

按照（111.9），空间的体积 $V$ 是与曲率半径 $a$ 的立方成比例的.因此 $\mathrm{d}V/V=3\,\mathrm{d}a/a=3\,\mathrm{d}\ln a$，于是我们可以写出
\begin{equation*}
  -\frac{\mathrm{d}\varepsilon}{\varepsilon+p}=3\,\mathrm{d}\ln a,
\end{equation*}
取积分则得
\begin{equation}
  3\ln a=-\int\frac{\mathrm{d}\varepsilon}{p+\varepsilon}+\mathrm{const}
\end{equation}
（积分的下限是常数）.

假如 $\varepsilon$ 与 $p$ 的关系（物态方程）是已知的，那么，由方程（112.6）就确定了 $\varepsilon$ 作为 $a$ 的函数.这时，从（112.5），我们可以确定 $\eta$ 如下：
\begin{equation}
  \eta=\pm\int\frac{\mathrm{d}a}
  {a\sqrt{\dfrac{8\pi k}{3c^{4}}(\varepsilon a^{2}-1)}}.
\end{equation}

方程（112.6）和（112.7）以普遍的形式解决了确定一个封闭的各向同性模型的度规的问题.

假如物质在空间是以不连续的宏观物体的形式分布的，那么，为了计算它所产生的引力场，我们可以将这些物体当做有一定质量的质点来处理，而完全不关注它们的内部构造.如果认为物体的速度较小（比光速 $c$ 小很多），我们可以简单地设 $\varepsilon=\mu c^{2}$，此处的 $\mu$ 是单位体积内的物体的质量之和.根据同样的道理，由这些物体构成的“气体”的压强比起 $\varepsilon$ 来是非常小的，因而可以略去不计（如我们所说，物体内部的压强与所考虑的问题无关）.至于空间中存在的辐射，其量相对地说也是很小的，因此，辐射能和辐射压也可以略去不计.

因此，为了用我们研究的模型来描述目前的宇宙状态，应该采用“尘埃状”物质的物态方程
\begin{equation*}
  \varepsilon=\mu c^{2},\quad p=0.
\end{equation*}

对（112.6）进行积分，就得到 $\mu a^{3}=\mathrm{const}$ 这个等式也可以直接写出，因为它不过说明了在整个空间内的物体的质量之和 $M$ 保持不变，这在我们研究的尘埃状物质的情形下是理所当然的\footnote{为了避免误解（读者在考虑到 §111 中提到的封闭宇宙的四维总动量等于零的评论时，有可能会产生这样的误解），我们强调，$M$ 是各个物体的质量之和，不考虑物体的引力相互作用.}.既然在闭合模型中空间的体积等于 $V=2\pi^{2}a^{3}$，那么 $\mathrm{const}=M/2\pi^{2}$.这样一来，
\begin{equation}
  \mu a^{3}=\mathrm{const}=\frac{M}{2\pi^{2}}.
\end{equation}

将（112.8）代入（112.7）并进行积分，我们得到
\begin{equation}
  a=a_{0}(1-\cos\eta),
\end{equation}
其中常数
\begin{equation*}
  a_{0}=\frac{2kM}{3\pi c^{2}}.
\end{equation*}
最后，对于 $t$ 和 $\eta$ 的关系，我们从（112.3）求得
\begin{equation}
  t=\frac{a_{0}}{c}(\eta-\sin\eta).
\end{equation}

方程（112.9），（112.10）决定以参数形式表示的函数 $a(t)$；函数 $a(t)$ 在 $t=0$（$\eta=0$）时刻从零开始增长，在 $t=\pi a_{0}/c$（$\eta=\pi$）时刻达到最大值 $a=2a_{0}$，然后在 $t=2\pi a_{0}/c$（$\eta=2\pi$）时刻又下降到零.

在 $\eta\ll1$ 时近似地有 $a=a_{0}\eta^{2}/2$，$t=a_{0}\eta^{3}/6c$，于是
\begin{equation}
  a\approx\left(\frac{9a_{0}c^{2}}{2}\right)^{1/3}t^{2/3}.
\end{equation}

在这种条件下物质的密度
\begin{equation}
  \mu=\frac{1}{6\pi kt^{2}}=\frac{8\times10^{5}}{t^{2}}
\end{equation}
（系数的数值按照以 $\mathrm{g\cdot cm^{-3}}$ 为单位的密度和以秒为单位的时间 $t$ 给出）.我们注意到，在这个范围内，函数 $\mu(t)$ 不依赖于参数 $a_{0}$，从这个角度来看，它具有普适的特点.

当 $a\to0$ 时密度 $\mu$ 变为无穷大.但是当 $\mu\to\infty$ 时，压强也会变得很大，因此为了研究上面所定的 $\eta$ 值附近的度规，我们必须考虑相反的极限情形，即尽可能大的压强的情形（对于给定的能量密度 $\varepsilon$ 来说），也就是说，用下面的物态方程来描述物质
\begin{equation*}
  p=\frac{\varepsilon}{3}
\end{equation*}
（参考§35第二个脚注）.于是从公式（112.6）我们得到
\begin{equation}
  \varepsilon a^{4}=\mathrm{const}\equiv\frac{3c^{4}a_{1}^{2}}{8\pi k}
\end{equation}
（式中 $a_{1}$ 是一个新的常数），在这以后，（112.7）和（112.3）导致下面的关系：
\begin{equation*}
  a=a_{1}\sin\eta,\quad t=\frac{a_{1}}{c}(1-\cos\eta).
\end{equation*}

既然这个解只对于很大的 $\varepsilon$（亦即很小的 $a$）才有意义，我们就假设 $\eta\ll1$.于是 $a\approx a_{1}\eta$，$t\approx a_{1}\eta^{2}/2c$，因此有
\begin{equation}
  a=\sqrt{2a_{1}ct}.
\end{equation}

在这种条件下
\begin{equation}
  \frac{\varepsilon}{c^{2}}=\frac{3}{32\pi kt^{2}}
  =\frac{4.5\times10^{5}}{t^{2}}
\end{equation}
（这个依赖关系还是不包含任何参数）.

于是，在 $t\to0$ 时仍然有 $a\to0$，因此数值 $t=0$ 确实是各向同性模型的时空度规的奇点（对于封闭模型中第二个 $a=0$ 点也是这样）.我们从（112.14）还能看到，在 $t$ 的符号发生改变时，$a(t)$ 变为虚数，而它的平方为负.这时式（112.2）中的所有四个分量 $g_{ik}$ 应该都是正的，而行列式 $g$ 也是正的.但是，这样的度规没有物理意义.这意味着，将该度规解析延拓到奇点之外是没有意义的.

\section{开放的各向同性模型}

用与上节完全相似的方法就可得到有负曲率的各向同性空间的相应的解（开放模型）.代替（112.2），现在我们有
\begin{equation}
  \mathrm{d}s^{2}=c^{2}\mathrm{d}t^{2}
  -a^{2}(t)\{\mathrm{d}\chi^{2}
  +\sinh^{2}\chi(\mathrm{d}\theta^{2}+\sin^{2}\theta\,\mathrm{d}\varphi^{2})\}.
\end{equation}

再次引入变量 $\eta$ 来代替 $t$，而 $\eta$ 与 $t$ 的关系是 $c\,\mathrm{d}t=a\,\mathrm{d}\eta$；这时，我们得到
\begin{equation}
  \mathrm{d}s^{2}=a^{2}(\eta)\{\mathrm{d}\eta^{2}-\mathrm{d}\chi^{2}
  -\sinh^{2}\chi(\mathrm{d}\theta^{2}+\sin^{2}\theta\,\mathrm{d}\varphi^{2})\}.
\end{equation}

在（112.4）式中，用 $\mathrm{i}\eta,\mathrm{i}\chi,\mathrm{i}a$ 分别代替 $\eta,\chi,a$，就能在形式上得到上式.因此，将相同的替换用于（112.5）和（112.6），我们也能直接得到场方程.这时，方程（112.6）保留它的原来形式：
\begin{equation}
  3\ln a=-\int\frac{\mathrm{d}\varepsilon}{\varepsilon+p}+\mathrm{const},
\end{equation}
而代替（112.5），我们有
\begin{equation}
  \frac{8\pi k}{c^{4}}\varepsilon=\frac{3}{a^{4}}(a^{2}-a^{2}).
\end{equation}

与此相应，代替（112.7），我们求得
\begin{equation}
  \eta=\pm\int\frac{\mathrm{d}a}
  {a\sqrt{\dfrac{8\pi k}{3c^{4}}\varepsilon a^{2}+1}}.
\end{equation}

对于尘埃状物质由此可以得到\footnote{我们指出，通过下面的变换
\[
r=A\mathrm{e}^{\eta}\sinh\chi,\qquad
c\tau=A\mathrm{e}^{\eta}\cosh\chi,\qquad
A\mathrm{e}^{\eta}=\sqrt{c^{2}\tau^{2}-r^{2}},\qquad
\tanh\chi=\frac{r}{c\tau},
\]
表达式（113.2）可以化为“共形伽利略”形式
\[
\mathrm{d}s^{2}=f(r,\tau)[c^{2}\mathrm{d}\tau^{2}-\mathrm{d}r^{2}-r^{2}(\mathrm{d}\theta^{2}+\sin^{2}\theta\,\mathrm{d}\varphi^{2})].
\]
具体地，在（113.6）的情形下我们得到（令 $A=a_{0}/2$）
\[
\mathrm{d}s^{2}=\left(1-\frac{a_{0}}{2\sqrt{c^{2}\tau^{2}-r^{2}}}\right)^{4}
\{c^{2}\mathrm{d}\tau^{2}-\mathrm{d}r^{2}-r^{2}(\mathrm{d}\theta^{2}+\sin^{2}\theta\,\mathrm{d}\varphi^{2})\}
\]
（V. A. Fock, 1955）. 在很大的值 $\sqrt{c^{2}\tau^{2}-r^{2}}$ 的情况下（对应于 $\eta\gg1$）这个度规趋向于伽利略度规，这种情况自然不出所料，因为曲率半径趋向无穷大.
在坐标 $r,\theta,\varphi,\tau$ 中，物质不是静止的，并且它的分布是不均匀的；在这种条件下物质的分布和运动围绕空间中的任何一点是中心对称的，该点被选作坐标 $\tau,\theta,\varphi$ 的原点.}：
\begin{equation}
  a=a_{0}(\cosh\eta-1),\quad t=\frac{a_{0}}{c}(\sinh\eta-\eta),
\end{equation}
\begin{equation}
  \mu a^{3}=\frac{3c^{2}}{4\pi k}a_{0}.
\end{equation}

公式（113.6）以参数形式决定了函数 $a(t)$.有别于封闭模型，这里曲率半径单调变化，在 $t=0$（$\eta=0$）时刻从零开始增加，在 $t\to\infty$（$\eta\to\infty$）时变为无穷大.相应地，物质的密度在 $t=0$ 时刻从无穷大开始单调减小（在 $\eta\ll1$ 时，这个减小的规律给出和封闭模型中相同的近似公式（112.12））.

对于很大的密度，解（113.6），（113.7）不能应用，必须再次转向 $p=\varepsilon/3$ 的情形.这种条件下再次得到关系式
\begin{equation}
  \varepsilon a^{4}=\mathrm{const}\equiv\frac{3c^{4}a_{1}^{2}}{8\pi k},
\end{equation}
而对于函数 $a(t)$，我们求得
\begin{equation*}
  a=a_{1}\sinh\eta,\quad t=\frac{a_{1}}{c}(\cosh\eta-1),
\end{equation*}
或者在 $\eta\ll1$ 时：
\begin{equation}
  a=\sqrt{2a_{1}ct}
\end{equation}
（以及以前得到的 $\varepsilon(t)$ 的公式（112.15））.这样一来，在开放模型中度规具有奇点（但和封闭模型的区别在于奇点只有一个）.

最后，所研究的解对应于空间的曲率半径无限大的极限情形是平直（欧几里得）空间模型.在这样的时空中的间隔 $\mathrm{d}s^{2}$ 可以写为
\begin{equation}
  \mathrm{d}s^{2}=c^{2}\mathrm{d}t^{2}
  -b^{2}(t)(\mathrm{d}x^{2}+\mathrm{d}y^{2}+\mathrm{d}z^{2})
\end{equation}
（我们选择了“笛卡儿”坐标 $x,y,z$ 作为空间坐标）.空间距离元中的时间因子显然不改变空间度规的欧几里得性，因为对于一个给定的 $t$，这个因子是一个常数；用简单的坐标变换就可以使之为1.用与上一节相似的计算，我们导出下面的方程：
\begin{equation*}
  \frac{8\pi k}{c^{2}}\varepsilon=\frac{3}{b^{2}}
  \left(\frac{\mathrm{d}b}{\mathrm{d}t}\right)^{2},\qquad
  3\ln b=-\int\frac{\mathrm{d}\varepsilon}{p+\varepsilon}+\mathrm{const}.
\end{equation*}

对于小压强的情形，我们求得
\begin{equation}
  \mu b^{3}=\mathrm{const},\quad b=\mathrm{const}\cdot t^{2/3}.
\end{equation}

对于小的 $t$，我们又必须考虑 $p=\varepsilon/3$ 的情形，对于这种情形，我们求得
\begin{equation}
  \varepsilon b^{4}=\mathrm{const},\quad b=\mathrm{const}\cdot\sqrt{t}.
\end{equation}

因此，在这种情形下，度规也有一个奇点（$t=0$）.

我们指出，所有得到的各向同性的解只有在物质的密度不为零的时候才能够存在；对于真空而言，爱因斯坦方程不具有这种类型的解\footnote{在 $\varepsilon=0$ 时从方程（113.5）我们就得到 $a=a_{0}\mathrm{e}^{\eta}=ct$（方程（112.7）由于虚数根一般会失去意义）. 但是度规
\[
\mathrm{d}s^{2}=c^{2}\mathrm{d}t^{2}-c^{2}t^{2}\{\mathrm{d}\chi^{2}+\sinh^{2}\chi(\mathrm{d}\theta^{2}+\sin^{2}\theta\,\mathrm{d}\varphi^{2})\}
\]
通过代换 $r=ct\sinh\chi$，$\tau=t\cosh\chi$，可化为形式
\[
\mathrm{d}s^{2}=c^{2}\mathrm{d}\tau^{2}-\mathrm{d}r^{2}-r^{2}(\mathrm{d}\theta^{2}+\sin^{2}\theta\,\mathrm{d}\varphi^{2}),
\]
即直接化为伽利略时空.}.同时我们提醒，从数学上的关系来看，它们是更加通用的解的类型的特殊情形，这些更加通用的解的类型包含三个物理上不同的空间坐标的任意函数（参见习题）.

\begin{xiti}

\noindent{\heiti 习题}\quad 假设在某个度规中空间的膨胀以“准均匀”的方式进行，就是说，所有的分量 $\gamma_{\alpha\beta}=-g_{\alpha\beta}$（在同步参考系中）以相同的规律趋于零，求这个度规在靠近奇点处的一般形式.空间充满了物态方程为 $p=\varepsilon/3$ 的物质（Е. М. Лифшиц, И. М. Халатников, 1960）.

解：在下面这样的形式中寻找靠近奇点处（$t=0$）的解：
\begin{equation}\tag{1}
  \gamma_{\alpha\beta}=ta_{\alpha\beta}+t^{2}b_{\alpha\beta}+\cdots,
\end{equation}
其中 $a_{\alpha\beta},b_{\alpha\beta}$ 是空间坐标的函数\footnote{弗里德曼解对应于函数 $a_{\alpha\beta}$ 的特殊选取，这种选取对应于恒定曲率的空间.}；下面设定 $c=1$.逆张量是
\begin{equation*}
  \gamma^{\alpha\beta}=\frac{1}{t}a^{\alpha\beta}-b^{\alpha\beta},
\end{equation*}
其中张量 $a^{\alpha\beta}$ 是 $a_{\alpha\beta}$ 的逆张量，而 $b^{\alpha\beta}=a^{\alpha\gamma}a^{\beta\delta}b_{\gamma\delta}$；下面所有升降指标的运算和协变微分均借助于不依赖于时间的度规 $a_{\alpha\beta}$ 来进行.

以需要的 $1/t$ 阶精度计算方程（97.11）和（97.12）的左边，我们得到
\begin{equation*}
  -\frac{3}{4t^{2}}+\frac{1}{2t}b=\frac{8\pi k}{3}\varepsilon(-4u_{0}^{2}+1),
  \qquad
  \frac{1}{2}(b_{;\alpha}-b_{\alpha;\beta}^{\beta})
  =-\frac{32\pi k}{3}\varepsilon u_{\alpha}u_{0}
\end{equation*}
（其中 $b=b_{\alpha}^{\alpha}$）.同时考虑恒等式
\begin{equation*}
  1=u_{i}u^{i}\approx u_{0}^{2}-\frac{1}{t}u_{\alpha}u_{\beta}a^{\alpha\beta},
\end{equation*}
我们得到
\begin{equation}\tag{2}
  8\pi k\varepsilon=\frac{3}{4t^{2}}-\frac{b}{2t},\qquad
  u_{\alpha}=\frac{t^{2}}{2}(b_{;\alpha}-b_{\alpha;\beta}^{\beta}).
\end{equation}

三维克里斯托夫符号，以及随之张量 $P_{\alpha\beta}$，在 $1/t$ 的一阶近似下不依赖于时间；在此情况下 $P_{\alpha\beta}$ 与只用度规 $a_{\alpha\beta}$ 所做的计算得到的表达式相同.考虑到这一点，我们得出在方程（97.13）中 $t^{-2}$ 量级的项相互抵消，而 $\sim1/t$ 的项给出
\begin{equation*}
  P_{\alpha}^{\beta}+\frac{3}{4}b_{\alpha}^{\beta}
  +\frac{5}{12}\delta_{\alpha}^{\beta}b=0,
\end{equation*}
由此
\begin{equation}\tag{3}
  b_{\alpha}^{\beta}=-\frac{4}{3}P_{\alpha}^{\beta}
  +\frac{5}{18}\delta_{\alpha}^{\beta}P
\end{equation}
（其中 $P=a^{\beta\gamma}P_{\beta\gamma}$）.鉴于恒等式
\begin{equation*}
  P_{\alpha;\beta}^{\beta}-\frac{1}{2}P_{;\alpha}=0
\end{equation*}
（参考（92.10））下面的关系式成立：
\begin{equation*}
  b_{\alpha;\beta}^{\beta}=\frac{7}{9}b_{;\alpha},
\end{equation*}
因此 $u_{\alpha}$ 可以写成如下形式
\begin{equation}\tag{4}
  u_{\alpha}=-\frac{t^{2}}{9}b_{;\alpha}.
\end{equation}

这样一来，所有的六个函数 $a_{\alpha\beta}$ 仍然是任意的，而展开式（1）中下一项的系数 $b_{\alpha\beta}$ 要按照它们确定.度规（1）中时间的选取完全取决于奇点处的条件 $t=0$；空间坐标还允许进行任意的不涉及时间的变换（例如，可以用它们将张量 $a_{\alpha\beta}$ 化为对角形式）.

因此得到的解总共包含三个“物理上不同”的任意函数.

我们指出，在这个解中空间度规是不均匀的和各向异性的，而在 $t\to0$ 时物质的密度分布趋向于均匀.三维速度 $\boldsymbol{v}$ 具有（在近似（4）中）等于零的旋度，而它的数值按照下面的规律趋近于零
\begin{equation*}
  v^{2}=v_{\alpha}v_{\beta}\gamma^{\alpha\beta}\sim t^{3}.
\end{equation*}

\end{xiti}

\section{红移}

所有我们研究过的解的基本特征是度规的非稳态性：空间的曲率半径是时间的函数.曲率半径的改变一般会导致空间内物体之间的所有距离的改变，这一点已经可以从空间距离元 $\mathrm{d}l$ 与 $a$ 成比例的情形看出.于是当 $a$ 增大时在这样的空间里物体彼此“退行”（在开放模型中 $a$ 的增大对应着 $\eta>0$，而在封闭模型中 $0<\eta<\pi$）.

假如有一个观察者坐在这些物体中的一个上面，在他看来，就好像所有其余物体都沿着视线方向离开他而去.这个“退行”速度（在给定的时刻 $t$）与物体之间的距离成比例.

这个预言有必要同一个基本的天文事实——星系谱线的红移——相比较.将这个移动解释为多普勒现象，我们就会得出星系“退行”的结论，就是说，在当今时代宇宙在膨胀\footnote{只有当物体的相互作用能量比它们“退行”时的动能小很多时，才能作出物体在 $a(t)$ 增大时彼此相互“退行”的结论；对分离得足够远的星系，这个条件总是满足的. 在相反的情形下，物体相互的距离基本上由它们的相互作用决定；因此，例如，这里所研究的效应实际上不应该影响到星系自身的尺度，对恒星尺度的影响就更小了.}.

我们来研究一条光线在各向同性空间中的传播.为了这个目的，最简便的方法是利用如下事实，即沿光信号传播的世界线，间隔 $\mathrm{d}s=0$.我们以光的出发点作为坐标 $\chi,\theta,\varphi$ 的原点.根据对称性的考虑，显而易见，光线是沿着“径向”传播的，即沿着 $\theta=\mathrm{const},\ \varphi=\mathrm{const}$ 的线传播.与此相应，我们在（112.4）或（113.2）中令 $\mathrm{d}\theta=\mathrm{d}\varphi=0$，就得到 $\mathrm{d}s^{2}=a^{2}(\mathrm{d}\eta^{2}-\mathrm{d}\chi^{2})$.令它等于零.我们得到 $\mathrm{d}\eta=\pm\mathrm{d}\chi$，或者，积分后得到
\begin{equation}
  \chi=\pm\eta+\mathrm{const}.
\end{equation}

对于从坐标原点发出的光线，$\eta$ 前取正号；而对于向原点传播的光线则取负号.这个形式的方程（114.1）既可应用到开放模型，也可以应用到封闭模型.借助于上节各公式，我们能够由此将光线传播的距离表示为时间的函数.

在开放模型中，一条光线从某一点出发，在传播过程中离开这一点越来越远.在封闭模型中，一条光线从起始点出发，最后能够到达空间的“对立极点”（这与 $\chi$ 从0变到 $\pi$ 相对应）；在以后的传播中，光线开始向起始点接近.光线绕着空间环行一周而又回到起始点的一个循环应该与 $\chi$ 从0变到 $2\pi$ 相对应.从（114.1）我们看到，这时 $\eta$ 也必须变化 $2\pi$，然而这是不可能的（光在与 $\eta=0$ 对应的时刻出发的情形除外）.因此一条光线在“绕空间”一周后不能回到出发点.

一条趋向观察点（坐标原点）传播的光线，与 $\eta$ 前面有负号的方程（114.1）相对应.假如光线到达这一点的时刻是 $t(\eta_{0})$，那么，当 $\eta=\eta_{0}$ 时，我们必定有 $\chi=0$，所以这样的光线的传播方程是
\begin{equation}
  \chi=\eta_{0}-\eta.
\end{equation}

由此可见，对于一个位于 $\chi=0$ 的点的观察者，只有那些从“距离”不超过 $\chi=\eta_{0}$ 的点出发的光线，才能在 $t(\eta_{0})$ 时刻到达观察者.

这个结果，对开放的和封闭的模型都很重要.我们看出，在时间 $t(\eta)$ 的每一时刻，在空间每一个给定点，不是整个空间都呈现在物理观察之下，只有与 $\chi\leqslant\eta$ 相对应的那一部分才可能被观察到.用数学的观点来看，空间的“可见区域”是四维时空被光锥所截的部分.这个部分无论对开放的模型或封闭的模型都是有限的（在开放模型中，被超曲面 $t=\mathrm{const}$ 所截的部分是无穷大的，这个超曲面对应于所有在同一时刻 $t$ 被观察到的点组成的空间）.在这个意义上，开放模型与封闭模型之间的差别并不如初见之时所想的那样深刻.

观察者在一定的时刻所观察到的区域离他愈远，与该区域相应的时刻就愈早.我们来看这样一个球面：这个球面是所有满足如下条件的点构成的几何图形，从这些点于 $t(\eta-\chi)$ 时刻发出的光于 $t(\eta)$ 时刻在原点被观察到.这个曲面的面积是 $4\pi a^{2}(\eta-\chi)\sin^{2}\chi$（在封闭模型中）或 $4\pi a^{2}(\eta-\chi)\sinh^{2}\chi$（在开放模型中）.当它远离观察者时，“可见球”的面积首先从零（当 $\chi=0$ 时）增加然后达到一个最大值，在此以后又减小，而当 $\chi=\eta$ 时又回到零（此处 $a(\eta-\chi)=a(0)=0$）.这就表明，光锥所割的截面不只是有限的，而且是封闭的.它好像是在与观察者“相对立”的点封闭；沿空间的任何方向进行观察，它都能被看到.在这一点 $\varepsilon\to\infty$，因此，物质在其所有的演化阶段原则上都是可观察的.

在开放模型中，所观察到的物质的总量等于
\begin{equation*}
  M_{\mathrm{obs}}=4\pi\int_{0}^{\eta}\mu a^{3}\sinh^{2}\chi\,\mathrm{d}\chi.
\end{equation*}
将（113.7）中的 $\mu a^{3}$ 代入，我们得到
\begin{equation}
  M_{\mathrm{obs}}=\frac{3c^{2}a_{0}}{2k}
  (\sinh\eta\cosh\eta-\eta).
\end{equation}

当 $\eta\to\infty$ 时，这个量无限制地增加.在封闭模型中，$M_{\mathrm{obs}}$ 的增加自然为总质量 $M$ 所限制.在这种情况下用类似的方式可以得到：
\begin{equation}
  M_{\mathrm{obs}}=\frac{M}{\pi}(\eta-\sin\eta\cos\eta).
\end{equation}

在 $\eta$ 从0增长到 $\pi$ 时这个量从0增长到 $M$；按照这个公式，接下来 $M_{\mathrm{obs}}$ 中远处的物体会被观察到两次（借助从两个方向“环绕空间”的光线）.

现在让我们考虑光在各向同性空间中传播时的频率变化.为此，我们首先指出下面的事实.设在空间的某一点有两个事件发生，这两个事件的时间间隔是 $\mathrm{d}t=a(\eta)\mathrm{d}\eta/c$.假如在这两个事件发生的两个时刻发出两个光信号，而在空间的另一点它们被观察到，那么，对应于光信号出发点量 $\eta$ 的变化 $\mathrm{d}\eta$，它们被观察到的两个时刻之间的时间间隔也对应相同的 $\mathrm{d}\eta$.直接从方程（114.1）可以推出这个结果.按照方程（114.1），$\eta$ 在光线从一点传播到另一点期间的变化仅与在这些点的坐标 $\chi$ 的差值有关.但是，既然在传播期间，曲率半径 $a$ 改变了，那么，发出信号的时刻与观察到它们的时刻的时间间隔 $t$ 就不同了，这些间隔之比等于相应的 $a$ 的值之比.

从上面所说的可以断定，例如，用世界时间 $t$ 来测量的光的振动周期也沿着光线改变，并与 $a$ 成比例.光的频率显然将与 $a$ 成反比.因此，在光线传播期间，沿着它的路径，下面的积为常数
\begin{equation}
  \omega a=\mathrm{const}.
\end{equation}

假设在 $t(\eta)$ 时刻我们观察从一个光源发出的光，这个光源所处的距离对应于坐标 $\chi$ 的一个确定值.按照（114.1）式，发射出这个光的时刻是 $t(\eta-\chi)$.假如 $\omega_{0}$ 是光在发射时刻的频率，那么，按照（114.5），我们所观察到的光的频率是
\begin{equation}
  \omega=\omega_{0}\frac{a(\eta-\chi)}{a(\eta)}.
\end{equation}

由于函数 $a(\eta)$ 的单调增加，我们有 $\omega<\omega_{0}$，即光频率将减小.这就是说，当我们观察向我们射来的光谱时，与在普通情形下同样物体的光谱相比较，所有它的光谱线必定移向红色的一边.红移现象在本质上是星系彼此“退行”的多普勒效应.\textsuperscript{*}{\renewcommand{\thefootnote}{*}\footnotetext{将星系红移解释为多普勒效应常用于红移较小的情况. 严格地说，这两种现象具有不同的物理起源. 多普勒效应是平直时空中的一种运动学效应，依赖于源和观测者的相对速度，与两者的距离无关；而星系红移主要源于宇宙膨胀，是大尺度时空弯曲的表现. 对于高红移天体，不能用狭义相对论的多普勒效应公式去推算其退行速度，而必须用基于广义相对论的（114.6）式.——中译注}}

利用移动了的频率与没有移动的频率之比 $\omega/\omega_{0}$ 测得的红移的大小，（对于给定的观察时刻）与被观察的光源所在处的距离有关（在（114.6）中，出现了光源的坐标 $\chi$）.当距离不太大时，我们可以将 $a(\eta-\chi)$ 展开为 $\chi$ 的幂级数，但只取前两项：
\begin{equation*}
  \frac{\omega}{\omega_{0}}=1-\chi\frac{a(\eta)}{a(\eta)}
\end{equation*}
（撇号表示对 $\eta$ 微分）.此外，我们指出，乘积 $\chi a(\eta)$ 在此恰恰是到被观察的光源的距离 $l$.事实上，“径向”线元等于 $\mathrm{d}l=a\,\mathrm{d}\chi$；在积分这个关系式时，就产生一个问题，即这个距离怎样通过物理观测来确定——由于这个原因在求这个距离时，我们必须在不同的时刻，在积分路线的不同点上取 $a$ 的值（$\eta=\mathrm{const}$ 的积分对应于同时观察沿路线上所有的点，这在物理上是不能实现的）.但是对于“小的”距离，我们可以略去 $a$ 沿积分路线的变化，而简单地写 $l=a\chi$，其中 $a$ 的值是在观察的时刻取的.

结果，我们求得频率改变的相对量如下（它通常用字母 $z$ 表示）：
\begin{equation}
  z=\frac{\omega_{0}-\omega}{\omega_{0}}=\frac{H}{c}\,l,
\end{equation}
此处我们引入了符号
\begin{equation}
  H=c\frac{a(\eta)}{a^{2}(\eta)}=\frac{1}{a}\frac{\mathrm{d}a}{\mathrm{d}t}
\end{equation}
作为所谓的哈勃常数.这个量对于给定的观察时刻与 $l$ 无关.因此光谱线的相对移动应当与到被观察光源的距离成比例.

通过将红移看做多普勒效应的结果，我们可以确定星系离开观察者的速度 $v$.写出 $z=v/c$，并与（114.7）相比较，我们得到
\begin{equation}
  v=Hl
\end{equation}
（直接计算导数 $v=\mathrm{d}(a\chi)/\mathrm{d}t$，也可以得到这个公式）.

天文学数据证实了规律（114.7），但是由于不能确定遥远星系的距离（这样的距离具有宇宙学尺度）这就给确定哈勃常数造成了麻烦.现今认可的 $H$ 的值为：
\begin{equation}
  \begin{gathered}
    H\approx0.8\times10^{-10}\,\mathrm{a}^{-1}
    =0.25\times10^{-17}\,\mathrm{s}^{-1},\\
    \frac{1}{H}\approx4\times10^{17}\,\mathrm{s}
    =1.3\times10^{10}\,\mathrm{a}.
  \end{gathered}
\end{equation}

这个值 $H$ 对应于“退行速度”在每个百万秒差距的距离上的增长为 $75\ \mathrm{km/s}$\footnote{还存在另一种估算，得出较小的 $H$，对应于“退行速度”在每个百万秒差距上的增长为 $55\ \mathrm{km/s}$，此时 $1/H=18\times10^{19}\,\mathrm{a}$.}.将 $\varepsilon=\mu c^{2}$ 和 $H=ca/a^{2}$ 代入（113.4）式，对于开放模型，我们得到关系式
\begin{equation}
  \frac{c^{2}}{a^{2}}=H^{2}-\frac{8\pi k}{3}\mu.
\end{equation}

将这个方程与等式
\begin{equation*}
  H=\frac{c\sinh\eta}{a_{0}(\cosh\eta-1)^{2}}
  =\frac{c}{a}\coth\frac{\eta}{2}
\end{equation*}
联立，我们得到
\begin{equation}
  \cosh\frac{\eta}{2}=H\sqrt{\frac{3}{8\pi k\mu}}.
\end{equation}

对于封闭模型，采用类似的方式，我们得到
\begin{equation}
  \frac{c^{2}}{a^{2}}=\frac{8\pi k}{3}\mu-H^{2},
\end{equation}
\begin{equation}
  \cos\frac{\eta}{2}=H\sqrt{\frac{3}{8\pi k\mu}}.
\end{equation}

比较（114.11）和（114.13），我们看出，空间曲率的正负取决于差量 $8\pi k\mu/3-H^{2}$ 的正负.对于 $\mu=\mu_{k}$，这个差量为零，其中
\begin{equation}
  \mu_{k}=\frac{3H^{2}}{8\pi k}.
\end{equation}

连同数值（114.10），我们得到 $\mu_{k}\approx1\times10^{-29}\ \mathrm{g/cm^{3}}$.根据现有的天文学知识，物质在空间的平均密度只能作很不准确的估计.基于星系的数量统计和它们的平均质量的估算，在现今取值约为 $3\times10^{-31}\ \mathrm{g/cm^{3}}$.这个值比 $\mu_{k}$ 小30倍，因而这个结果对开放模型有利.然而，暂不说这个值本身不太充分的可信度，应该意识到，它并没有考虑星系之间可能存在的暗气体，考虑这些暗气体能够大大提高物质的平均密度.\textsuperscript{*}{\renewcommand{\thefootnote}{*}\footnotetext{新近的观测表明，宇宙的总能量密度非常接近 $\mu_{k}$，其中，包括暗物质在内的物质约占 30\%，暗能量约占 70\%.——中译注}}

让我们在这里指出一些不等式，在给定 $H$ 之值的情况下可以得到这些不等式.对于开放模型，我们有 $H=c\sinh\eta/[a_{0}(\cosh\eta-1)^{2}]$，于是
\begin{equation*}
  t=\frac{a_{0}}{c}(\sinh\eta-\eta)
  =\frac{\sinh\eta(\sinh\eta-\eta)}{H(\cosh\eta-1)^{2}}.
\end{equation*}
既然 $0<\eta<\infty$，那么我们应当有
\begin{equation}
  \frac{2}{3H}<t<\frac{1}{H}.
\end{equation}

同理，对于封闭模型，我们得到
\begin{equation*}
  t=\frac{\sin\eta(\eta-\sin\eta)}{H(1-\cos\eta)^{2}}.
\end{equation*}
$a(\eta)$ 的增加对应于区间 $0<\eta<\pi$；因此我们得到
\begin{equation}
  0<t<\frac{2}{3H}.
\end{equation}

下面，我们决定光从光源到达观察者时的强度 $I$，光源到观察者的距离与坐标 $\chi$ 的一定值相对应.在观察点光的能流密度与球的面积成反比，该球面以光源所在点为球心并经过被考察点；在负曲率的空间中，球的面积等于 $4\pi a^{2}\sinh^{2}\chi$.此外，光源在时间间隔 $\mathrm{d}t=a(\eta-\chi)\mathrm{d}\eta/c$ 内所射出的光会在时间间隔 $a(\eta)\mathrm{d}t/a(\eta-\chi)=a(\eta)\mathrm{d}\eta/c$ 内达到观察者所在之点.因为光的强度被定义为每单位时间的光能流量，所以在 $I$ 内就出现了一个因子 $a(\eta-\chi)/a(\eta)$.最后，一个波包的能量与它的频率成比例（参考（53.9））；既然在光传播期间，频率按照规律（114.5）而改变，因而这导致 $I$ 中又出现一个因子 $a(\eta-\chi)/a(\eta)$.结果，我们得到强度的公式如下：
\begin{equation}
  I=\mathrm{const}\,\frac{a^{2}(\eta-\chi)}
  {a^{4}(\eta)\sinh^{2}\chi}.
\end{equation}

对于封闭模型，我们同样地得到
\begin{equation}
  I=\mathrm{const}\,\frac{a^{2}(\eta-\chi)}
  {a^{4}(\eta)\sin^{2}\chi}.
\end{equation}

这两个公式决定被观察对象的视亮度与其距离的关系（在绝对亮度一定的情况下）.对于小的 $\chi$，我们可以令 $a(\eta-\chi)\approx a(\eta)$，这样，我们就有 $I\sim1/a^{2}(\eta)\chi^{2}=1/l^{2}$，即是说，我们得到了光强度与距离的平方成反比的普通定律.

最后，让我们考虑所谓物体的本动问题.当说到物质的密度和运动时，我们总是理解为平均密度和平均运动；特别是，在我们一直使用的参考系中，平均运动的速度为零.各物体的实际速度围绕这个平均值呈现出一定的涨落.在时间的进程中，物体的本动速度是变化的.为了决定这个变化的规律，我们来考虑一个自由运动的物体，并且选择轨道上的任意一点为坐标原点，那么，轨道将是径向线：$\theta=\mathrm{const},\ \varphi=\mathrm{const}$.将 $g^{ik}$ 的值代入以后，哈密顿-雅可比方程（87.6）取以下的形式：
\begin{equation}
  \left(\frac{\partial S}{\partial\chi}\right)^{2}
  -\left(\frac{\partial S}{\partial\eta}\right)^{2}
  +m^{2}c^{2}a^{2}(\eta)=0.
\end{equation}

因为 $\chi$ 没有出现在这个方程的系数内（也就是说 $\chi$ 是一个循环坐标），那么，守恒定律 $\partial S/\partial\chi=\mathrm{const}$ 将是有效的.根据一般的定义，运动物体的动量 $p=\partial S/\partial l=\partial S/a\,\partial\chi$.因此，对于运动的物体，下面的乘积是常数：
\begin{equation}
  pa=\mathrm{const}.
\end{equation}

按照下式引入物体本动的速度 $v$
\begin{equation*}
  p=\frac{mv}{\sqrt{1-\dfrac{v^{2}}{c^{2}}}},
\end{equation*}
我们得到
\begin{equation}
  \frac{va}{\sqrt{1-\dfrac{v^{2}}{c^{2}}}}=\mathrm{const}.
\end{equation}

这些关系决定了速度随时间变化的规律.随着 $a$ 的增加，速度 $v$ 单调地减小.

\begin{xiti}

\noindent{\heiti 习题1}\quad 求出星系的视亮度作为红移函数的展开式的头两项，星系绝对亮度随时间按照指数规律 $I_{\mathrm{abs}}=\mathrm{const}\cdot\mathrm{e}^{\alpha t}$ 变化（H. Robertson, 1955）.

解：在“时刻”$\eta$ 观察到的星系的视亮度与距离 $\chi$ 的关系由下面公式给出（对于封闭模型）
\begin{equation*}
  I=\mathrm{const}\cdot\mathrm{e}^{\alpha[t(\eta-\chi)-t(\eta)]}
  \frac{a^{2}(\eta-\chi)}{a^{4}(\eta)\sin^{2}\chi}.
\end{equation*}
按照（114.7）定义的红移
\begin{equation*}
  z=\frac{\omega_{0}-\omega}{\omega}
  =\frac{a(\eta)-a(\eta-\chi)}{a(\eta-\chi)}.
\end{equation*}
将 $I$ 和 $z$ 按 $\chi$ 的幂展开（连同（112.9），（112.10）中的函数 $a(\eta)$ 和 $t(\eta)$）然后从得到的表达式中消去 $\chi$，结果我们求得
\begin{equation*}
  I=\mathrm{const}\cdot\frac{1}{z^{2}}
  \left[1-\left(1-\frac{q}{2}+\frac{\alpha}{H}\right)z\right],
\end{equation*}
其中引入符号
\begin{equation*}
  q=\frac{2}{1+\cos\eta}=\frac{\mu}{\mu_{k}}>1.
\end{equation*}
对于开放模型得到同样的公式，但其中
\begin{equation*}
  q=\frac{2}{1+\cosh\eta}=\frac{\mu}{\mu_{k}}<1.
\end{equation*}

\noindent{\heiti 习题2}\quad 求给定半径的球面内星系数量作为球面边界上红移函数的展开式的头几项（假设星系的空间分布是均匀的）.

解：位于 $\leqslant\chi$ 的“距离”范围内的星系的数量 $N$ 为（在封闭模型中）
\begin{equation*}
  N=\mathrm{const}\cdot\int_{0}^{\chi}\sin^{2}\chi\,\mathrm{d}\chi
  \approx\mathrm{const}\cdot\chi^{3}.
\end{equation*}
向这里代入函数 $\chi(z)$ 的展开式的头两项，我们得到：
\begin{equation*}
  N=\mathrm{const}\cdot z^{3}
  \left[1-\frac{3}{4}(2+q)z\right].
\end{equation*}
这个形式的公式对于开放模型也成立.

\end{xiti}

\section{各向同性宇宙的引力稳定性}

我们来研究关于各向同性模型中微小扰动的演化特征，也就是说，关于它的引力稳定性问题（E. M. Lifshitz, 1946）.我们仅限于研究在相对不大的空间区域内的扰动——该空间区域的线度和半径 $a$ 相比是小量\footnote{对这个问题更详细的叙述，包括在线度与 $a$ 相当的区域内的扰动的研究可参考 E. M. Lifshitz//УФН. 1963. Т. 80. С. 411; Adv. in Phys. 1963. V. 12. P. 208.}.

在每一个这样的区域内，空间度规在一阶近似中可以取作欧几里得度规，就是说，度规（111.8）或（111.12）可替换为度规
\begin{equation}
  \mathrm{d}l^{2}=a^{2}(\eta)
  (\mathrm{d}x^{2}+\mathrm{d}y^{2}+\mathrm{d}z^{2}),
\end{equation}
其中 $x,y,z$ 是以 $a$ 为度量单位的笛卡儿坐标.我们还将像从前那样使用变量 $\eta$ 作为时间坐标.

不失一般性，我们将像从前那样在同步参考系中描述被扰动的场，就是说给度规张量的变分 $\delta g_{ik}$ 加上条件 $\delta g_{00}=\delta g_{0\alpha}=0$.在这些条件下对恒等式 $g_{ik}u^{i}u^{k}=1$ 进行变分（应注意到，物质四维速度分量的无扰动的值为 $u^{0}=1/a,\ u^{\alpha}=0$\footnote{在本节中，不再用附加上标 $(0)$ 来表示未扰动的物理量.}），我们得到 $g_{00}u^{0}\delta u^{0}=0$，由此 $\delta u^{0}=0$.扰动 $\delta\boldsymbol{u}^{\alpha}$ 一般说来是不为零的，因此这个参考系已经不是共动参考系.

我们利用 $h_{\alpha\beta}=\delta\gamma_{\alpha\beta}=-\delta g_{\alpha\beta}$ 表示空间度规张量的扰动，于是 $\delta\gamma^{\alpha\beta}=-h^{\alpha\beta}$，这里借助未扰动度规 $\gamma_{\alpha\beta}$ 来提升 $h_{\alpha\beta}$ 的指标.

在线性近似下引力场的微小扰动满足方程
\begin{equation}
  \delta R_{i}^{k}-\frac{1}{2}\delta_{i}^{k}\delta R
  =\frac{8\pi k}{c^{4}}\delta T_{i}^{k}.
\end{equation}

在同步参考系中能量动量张量（94.9）的分量的变分等于
\begin{equation}
  \delta T_{\alpha}^{\beta}=-\delta_{\alpha}^{\beta}\delta p,\quad
  \delta T_{0}^{\alpha}=a(p+\varepsilon)\delta u^{\alpha},\quad
  \delta T_{0}^{0}=\delta\varepsilon.
\end{equation}

鉴于 $\delta\varepsilon$ 和 $\delta p$ 都是小量，可以写出 $\delta p=\dfrac{\mathrm{d}p}{\mathrm{d}\varepsilon}\delta\varepsilon$，然后我们得到关系式
\begin{equation}
  \delta T_{\alpha}^{\beta}
  =-\delta_{\alpha}^{\beta}\frac{\mathrm{d}p}{\mathrm{d}\varepsilon}\delta T_{0}^{0}.
\end{equation}

针对 $\delta R_{i}^{k}$ 的公式可以通过对（97.10）进行变分来获得.由于未扰动的度规张量 $\gamma_{\alpha\beta}=a^{2}\delta_{\alpha\beta}$，那么未扰动值是
\begin{equation*}
  \varkappa_{\alpha\beta}=\frac{2\dot{a}}{a}\gamma_{\alpha\beta}
  =\frac{2a}{a^{2}}\gamma_{\alpha\beta},\qquad
  \varkappa_{\alpha}^{\beta}=\frac{2a}{a^{2}}\delta_{\alpha}^{\beta},
\end{equation*}
其中符号上方的点表示对 $ct$ 微分，而撇表示对 $\eta$ 微分.量 $\varkappa_{\alpha\beta}$ 和 $\varkappa_{\alpha}^{\beta}=\varkappa_{\alpha\gamma}\gamma^{\gamma\beta}$ 的扰动如下：
\begin{equation*}
  \delta\varkappa_{\alpha\beta}=\dot{h}_{\alpha\beta}=\frac{1}{a}h_{\alpha\beta},
  \qquad
  \delta\varkappa_{\alpha}^{\beta}
  =-h^{\beta\gamma}\varkappa_{\alpha\gamma}
  +\gamma^{\beta\gamma}\dot{h}_{\alpha\gamma}
  =\dot{h}_{\alpha}^{\beta}=\frac{1}{a}h_{\alpha}^{\beta},
\end{equation*}
其中 $h_{\alpha}^{\beta}=\gamma^{\beta\gamma}h_{\alpha\gamma}$.对于欧几里得度规（115.1），三维张量 $P_{\alpha}^{\beta}$ 的未扰动值等于零.$\delta P_{\alpha}^{\beta}$ 的变分按照公式（108.3），（108.4）计算；显然，正如 $\delta R_{ik}$ 可用 $\delta g_{ik}$ 表达，$\delta P_{\alpha}^{\beta}$ 可以用 $\delta\gamma_{\alpha\beta}$ 来表达，所有的张量运算都在带有度规（115.1）的三维空间里进行；鉴于这个度规是欧几里得的，所有的协变微分可简化为对坐标 $x^{\alpha}$ 的普通微分（对于逆变微分仍需再除以 $a^{2}$）.考虑到所有这些情况（并且从对 $t$ 的导数全面过渡到对 $\eta$ 的导数），在简单的计算之后我们得到：
\begin{equation}
  \begin{aligned}
    &\delta R_{\alpha}^{\beta}
    =-\frac{1}{2a^{2}}\bigl(h_{\alpha,\gamma}^{\gamma,\beta}
    +h_{\gamma,\alpha}^{\beta,\gamma}
    -h_{\alpha,\gamma}^{\beta,\gamma}
    -h_{,\alpha}^{\beta}\bigr)
    -\frac{1}{2a^{2}}h_{\alpha}^{\beta}
    -\frac{a}{a^{3}}h_{\alpha}^{\beta}
    -\frac{a}{2a^{3}}h\delta_{\alpha}^{\beta},\\
    &\delta R_{0}^{0}=-\frac{1}{2a^{2}}h
    -\frac{a}{2a^{3}}h,\\
    &\delta R_{0}^{\alpha}
    =\frac{1}{2a^{2}}\bigl(h^{,\alpha}-h_{\beta}^{\alpha,\beta}\bigr),\\
    &(h\equiv h_{\alpha}^{\alpha}).
  \end{aligned}
\end{equation}

这里逗号后面的指标无论在上方还是在下方，都表示对坐标 $x^{\alpha}$ 的简单微分（我们继续在上方和下方书写指标，只是为了保持表示方法的统一性）.

把通过 $\delta R_{i}^{k}$ 表达的分量 $\delta T_{i}^{k}$ 按照（115.2）式代入（115.4）式，我们将得到关于扰动 $h_{\alpha}^{\beta}$ 的最终方程.就这些方程而言，选择在 $\alpha\neq\beta$ 情形下的（115.4），以及按指标 $\alpha,\beta$ 进行缩并而得的方程是便利的，它们是：
\begin{equation}
  \begin{aligned}
    &\bigl(h_{\alpha,\gamma}^{\gamma,\beta}
    +h_{\gamma,\alpha}^{\beta,\gamma}
    -h_{,\alpha}^{\beta}
    -h_{\alpha,\gamma}^{\beta,\gamma}\bigr)
    +h_{\alpha}^{\beta}
    +2\frac{a}{a}h_{\alpha}^{\beta}=0,\qquad\alpha\neq\beta,\\
    &\frac{1}{2}\bigl(h_{\gamma,\delta}^{\delta,\gamma}
    -h_{,\gamma}^{\gamma}\bigr)
    \left(1+3\frac{\mathrm{d}p}{\mathrm{d}\varepsilon}\right)
    +h+h\frac{a}{a}
    \left(2+3\frac{\mathrm{d}p}{\mathrm{d}\varepsilon}\right)=0.
  \end{aligned}
\end{equation}

物质密度和速度的扰动可以按照已知的 $h_{\alpha}^{\beta}$，借助于公式（115.2），（115.3）来确定.因此，对于密度的相对变化，我们有：
\begin{equation}
  \frac{\delta\varepsilon}{\varepsilon}
  =\frac{c^{4}}{8\pi k\varepsilon}
  \left(\delta R_{0}^{0}-\frac{1}{2}\delta R\right)
  =\frac{c^{4}}{16\pi k\varepsilon a^{2}}
  \left(h_{\alpha,\beta}^{\beta,\alpha}
  -h_{,\alpha}^{,\alpha}
  +\frac{2a}{a}h\right).
\end{equation}

在方程（115.6）的解当中存在这样的一些解，它们可以通过简单的参考系变换（不破坏其同步性）而被消去，因此它们不是度规的实际物理变化.这些解的形式可以借助于§97节习题3中得到的公式（1）和（2）提前确立.将未扰动的值 $\gamma_{\alpha\beta}=a^{2}\delta_{\alpha\beta}$ 代入其中，我们将得到下列针对度规的虚拟扰动的表达式：
\begin{equation}
  h_{\alpha}^{\beta}
  =f_{0,\alpha}{}^{,\beta}\int\frac{\mathrm{d}\eta}{a}
  +\frac{a}{a^{2}}f_{0}\delta_{\alpha}^{\beta}
  +(f_{\alpha}{}^{,\beta}+f^{\beta}{}_{,\alpha}),
\end{equation}
其中 $f_{0},f_{\alpha}$ 是坐标 $x,y,z$ 的任意（小量的）函数.

由于度规在我们考察的不大的空间区域内设定为欧几里得的，那么任意的扰动在每一个这样的区域可以分解成平面波.将 $x,y,z$ 作为以 $a$ 为度量单位的笛卡儿坐标，我们可以将平面波的周期性空间因子写成 $\mathrm{e}^{\mathrm{i}\boldsymbol{n}\cdot\boldsymbol{r}}$ 的形式，其中 $\boldsymbol{n}$ 为无量纲矢量，表示以 $1/a$ 为度量单位的波矢（波矢 $\boldsymbol{k}=\boldsymbol{n}/a$）.如果在尺度 $\sim l$ 的空间区域内有扰动，那么在它的展开式中会出现波长 $\lambda=2\pi a/n\sim l$ 的波.对于局限在尺度 $l\ll a$ 的区域内的扰动，我们相应地假定数 $n$ 足够大（$n\gg2\pi$）.

引力扰动可以分为三个类型.这个分类可归结为确定平面波的可能种类，对称张量 $h_{\alpha\beta}$ 可以用这些不同种类的平面波表示出来.这样，我们得到下面的分类：

\subsection*{1. 借助于标量函数}
\begin{equation}
  Q=\mathrm{e}^{\mathrm{i}\boldsymbol{n}\cdot\boldsymbol{r}},
\end{equation}
可以构造矢量 $\boldsymbol{P}=\boldsymbol{n}Q$ 和张量\footnote{我们书写普通的笛卡儿矢量 $\boldsymbol{n}$ 的分量时也区分上下指标，这样做仅仅是为了保持表示方法的统一性.}
\begin{equation}
  Q_{\alpha}^{\beta}=\frac{1}{3}\delta_{\alpha}^{\beta}Q,\qquad
  P_{\alpha}^{\beta}
  =\left(\frac{1}{3}\delta_{\alpha}^{\beta}
  -\frac{n_{\alpha}n^{\beta}}{n^{2}}\right)Q.
\end{equation}

这些平面波对应于这样的扰动，其中物质的速度和密度与引力场一起经历着变化，就是说，我们在处理伴随着物质聚集和疏松的扰动.在这种情况下扰动 $h_{\alpha}^{\beta}$ 通过张量 $Q_{\alpha}^{\beta}$ 和 $P_{\alpha}^{\beta}$ 表达，速度的扰动通过矢量 $\boldsymbol{P}$，而密度的扰动通过标量 $Q$ 表达.

\subsection*{2. 借助于矢量横波}
\begin{equation}
  \boldsymbol{S}=\boldsymbol{s}\,\mathrm{e}^{\mathrm{i}\boldsymbol{n}\cdot\boldsymbol{r}},
  \qquad
  \boldsymbol{s}\cdot\boldsymbol{n}=0,
\end{equation}
可以构造张量 $(n^{\beta}S_{\alpha}+n_{\alpha}S^{\beta})$；由于 $\boldsymbol{n}\cdot\boldsymbol{S}=0$，对应的标量不存在.这些波对应于这样的扰动，其中速度与引力场一起经历着变化，但不包括物质的密度；它们可以被称作旋转扰动.

\subsection*{3. 张量横波}
\begin{equation}
  G_{\alpha}^{\beta}=g_{\alpha}^{\beta}\mathrm{e}^{\mathrm{i}\boldsymbol{n}\cdot\boldsymbol{r}},
  \qquad
  g_{\alpha}^{\beta}n_{\beta}=0.
\end{equation}

利用张量横波既无法构造矢量，也无法构造标量.这些波对应于引力场这样的扰动，其中物质不运动，并且在空间中均匀分布.换句话说，这是各向同性宇宙中的引力波.

最有意义的是第一种类型的扰动.令
\begin{equation}
  h_{\alpha}^{\beta}=\lambda(\eta)P_{\alpha}^{\beta}
  +\mu(\eta)Q_{\alpha}^{\beta},\qquad h=\mu Q.
\end{equation}

从（115.7）中可以得到密度的相对变化：
\begin{equation}
  \frac{\delta\varepsilon}{\varepsilon}
  =\frac{c^{4}}{24\pi k\varepsilon a^{2}}
  \left[n^{2}(\lambda+\mu)+\frac{3a}{a}\mu\right]Q.
\end{equation}

将（115.13）代入（115.6），得到确定函数 $\lambda$ 和 $\mu$ 的方程
\begin{equation}
  \begin{aligned}
    &\lambda+2\frac{a}{a}\lambda
    -\frac{n^{2}}{3}(\lambda+\mu)=0,\\
    &\mu+\mu\frac{a}{a}
    \left(2+3\frac{\mathrm{d}p}{\mathrm{d}\varepsilon}\right)
    +\frac{n^{2}}{3}(\lambda+\mu)
    \left(1+3\frac{\mathrm{d}p}{\mathrm{d}\varepsilon}\right)=0.
  \end{aligned}
\end{equation}

这些方程具有下列两个特殊积分，对应于可以通过变换参考系而消除的度规的虚拟变化：
\begin{equation}
  \lambda=-\mu=\mathrm{const},
\end{equation}
\begin{equation}
  \lambda=-n^{2}\int\frac{\mathrm{d}\eta}{a},\qquad
  \mu=n^{2}\int\frac{\mathrm{d}\eta}{a}-\frac{3a}{a^{2}}
\end{equation}
（其中第一个是从（115.8）式通过选取 $f_{0}=0,\ f_{\alpha}=P_{\alpha}$ 得出，第二个是通过选取 $f_{0}=Q,\ f_{\alpha}=0$ 得出）.

在宇宙膨胀的早期，那时的物质用物态方程 $p=\varepsilon/3$ 描述，我们有 $a\approx a_{1}\eta$，$\eta\ll1$（在开放模型与封闭模型中都是这样）.方程（115.15）取如下形式：
\begin{equation}
  \lambda+\frac{2}{\eta}\lambda
  -\frac{n^{2}}{3}(\lambda+\mu)=0,\qquad
  \mu+\frac{3}{\eta}\mu
  +\frac{2n^{2}}{3}(\lambda+\mu)=0.
\end{equation}

针对两种极限情况对这些方程分别进行研究是方便的，这两种极限情况取决于两个较大的量 $n$ 和 $1/\eta$ 之间的相互关系.

首先我们假设 $n$ 不是非常大（或者 $\eta$ 足够小），于是 $n\eta\ll1$.在这种情况下，以能够使方程（115.18）成立所需要的精度，由这些方程我们可以得出：
\begin{equation*}
  \lambda=\frac{3C_{1}}{\eta}
  +C_{2}\left(1+\frac{n^{2}}{9}\eta^{2}\right),\qquad
  \mu=-\frac{2n^{2}}{3}C_{1}\eta
  +C_{2}\left(1-\frac{n^{2}}{6}\eta^{2}\right),
\end{equation*}
其中 $C_{1},C_{2}$ 为常数；从这里排除了形式为（115.16）和（115.17）的解（在当前情况下，这些解一个有 $\lambda=-\mu=\mathrm{const}$，一个有 $\lambda+\mu\sim1/\eta^{2}$）.同时按照（115.14）和（112.15）计算出 $\delta\varepsilon/\varepsilon$，我们得到下列关于度规和密度扰动的表达式：
\begin{equation}
  \begin{gathered}
    h_{\alpha}^{\beta}=\frac{3C_{1}}{\eta}P_{\alpha}^{\beta}
    +C_{2}(Q_{\alpha}^{\beta}+P_{\alpha}^{\beta}),\\
    \frac{\delta\varepsilon}{\varepsilon}
    =\frac{n^{2}}{9}(C_{1}\eta+C_{2}\eta^{2})Q
    \quad\text{当}\quad p=\frac{\varepsilon}{3},\quad
    \eta\ll\frac{1}{n}.
  \end{gathered}
\end{equation}

常数 $C_{1},C_{2}$ 必须满足特定的条件，这些条件应能反映出扰动在其产生时刻 $\eta_{0}$ 是很小的：应该有 $h_{\alpha}^{\beta}\ll1$（由此 $\lambda\ll1,\ \mu\ll1$）和 $\delta\varepsilon/\varepsilon\ll1$.将这些条件用于（115.19）可导出不等式 $C_{1}\ll\eta_{0},\ C_{2}\ll1$.

在（115.19）式中有一些项，它们在膨胀的宇宙中以半径 $a=a_{1}\eta$ 的不同次幂增长.但是这个增长并不会导致扰动能够成为大的量：如果我们应用公式（115.19）精确到 $\eta\sim1/n$ 的数量级，那么可以看到，（基于上面得到的关于 $C_{1},C_{2}$ 的不等式）扰动即使在这些公式的应用上限仍然是小量.

现在假设 $n$ 很大，使得有 $n\eta\gg1$.在这个条件下求解方程（115.18），我们得到 $\lambda$ 和 $\mu$ 中占主导地位的项是：\footnote{指数项前面的因数 $1/\eta^{2}$ 是按 $1/n\eta$ 的幂展开式的第一项. 在当前情况下，为了确定它，应该同时考虑展开式的头两项（方程（115.18）的精度是允许这样做的）.}
\begin{equation*}
  \lambda=-\frac{\mu}{2}
  =\mathrm{const}\cdot\frac{1}{\eta^{2}}
  \mathrm{e}^{\mathrm{i}n\eta/\sqrt{3}}.
\end{equation*}

由此我们得到度规和密度的扰动：
\begin{equation}
  \begin{gathered}
    h_{\alpha}^{\beta}
    =\frac{C}{n^{2}\eta^{2}}
    \bigl(P_{\alpha}^{\beta}-2Q_{\alpha}^{\beta}\bigr)
    \mathrm{e}^{\mathrm{i}n\eta/\sqrt{3}},\qquad
    \frac{\delta\varepsilon}{\varepsilon}
    =-\frac{C}{9}Q\,\mathrm{e}^{\mathrm{i}n\eta/\sqrt{3}}\\
    \text{对于条件}\ p=\frac{\varepsilon}{3},\quad
    \frac{1}{n}\ll\eta\ll1,
  \end{gathered}
\end{equation}
其中 $C$ 为复数常数，满足条件 $|C|\ll1$.在这些表达式中周期性因子的出现是十分自然的.在 $n$ 为大数的条件下我们是在处理这样的扰动，其空间的周期性取决于大的波矢 $k=n/a$.这样的扰动应该像声波那样传播，其速度为
\begin{equation*}
  u=\sqrt{\frac{\mathrm{d}p}{\mathrm{d}(\varepsilon/c^{2})}}
  =\frac{c}{\sqrt{3}}.
\end{equation*}

相应地，相位的时间部分就像几何声学里规定的那样，通过大积分 $\int ku\,\mathrm{d}t=n\eta/\sqrt{3}$ 来确定.正如我们所见，密度相对变化的振幅保持恒定，而度规扰动的振幅在宇宙膨胀的情况下按照 $a^{-2}$ 减小\footnote{容易验证（在 $p=\varepsilon/3$ 的条件下）$n\eta\sim L/\lambda$，其中 $L\sim u/\sqrt{k\varepsilon/c^{2}}$. 理所当然，特征长度 $L$（它决定波长 $\lambda\ll a$ 的扰动行为）仅仅包含“流体力学”量——物质密度 $\varepsilon/c^{2}$ 和声速 $u$（还有引力常数 $k$）. 我们指出，在 $\lambda\gg L$ 的条件下扰动会增长（见（115.19））.}.

接下来，我们研究膨胀的更晚阶段，那时物质已经稀薄到可以忽略掉其压强的程度（$p=0$）.在这种条件下，我们此处仅局限于 $\eta$ 很微小的情况，在这些微小的 $\eta$ 对应的膨胀阶段内，半径 $a$ 和它现今的值相比还很小，但尽管如此物质已经足够稀薄了.

在 $p=0$ 和 $\eta\ll1$ 的条件下我们有 $a\approx a_{0}\eta^{2}/2$，并且方程（115.15）具有如下形式
\begin{equation*}
  \lambda+\frac{4}{\eta}\lambda
  -\frac{n^{2}}{3}(\lambda+\mu)=0,
\end{equation*}
\begin{equation*}
  \mu+\frac{4}{\eta}\mu
  +\frac{n^{2}}{3}(\lambda+\mu)=0.
\end{equation*}

这些方程的解是：
\begin{equation*}
  \lambda+\mu=2C_{1},\qquad
  \lambda-\mu=2n^{2}\left(\frac{C_{1}\eta^{2}}{15}
  +\frac{2C_{2}}{\eta^{3}}\right).
\end{equation*}

也可算出 $\delta\varepsilon/\varepsilon$（利用（115.14）和（112.12）），我们求得
\begin{equation}
  \begin{gathered}
    h_{\alpha}^{\beta}
    =C_{1}(P_{\alpha}^{\beta}+Q_{\alpha}^{\beta})
    +\frac{2n^{2}C_{2}}{\eta^{3}}
    (P_{\alpha}^{\beta}-Q_{\alpha}^{\beta})
    \quad\text{当}\quad\eta\ll\frac{1}{n},\\
    h_{\alpha}^{\beta}
    =\frac{C_{1}}{15}n^{2}\eta^{2}(P_{\alpha}^{\beta}-Q_{\alpha}^{\beta})
    +\frac{2n^{2}C_{2}}{\eta^{3}}
    (P_{\alpha}^{\beta}-Q_{\alpha}^{\beta})
    \quad\text{当}\quad\frac{1}{n}\ll\eta\ll1,\\
    \frac{\delta\varepsilon}{\varepsilon}
    =\left(\frac{C_{1}n^{2}\eta^{2}}{30}
    +\frac{C_{2}n^{2}}{\eta^{3}}\right)Q.
  \end{gathered}
\end{equation}

我们看到，$\delta\varepsilon/\varepsilon$ 包含正比于 $a$ 的增长项\footnote{考虑微小压强 $p(\varepsilon)$ 的更细致的分析表明，满足条件 $u\eta n/c\ll1$ 时才有可能忽略压强（其中 $u=c\sqrt{\mathrm{d}p/\mathrm{d}\varepsilon}$ 是微小的声速）；容易验证，这个条件和条件 $\lambda\gg L$ 是相同的. 这样一来，只要 $\lambda\gg L$，扰动总是会增长.}.但是如果 $n\eta\ll1$，那么按照条件 $C_{1}\ll1$，即使当 $\eta\sim1/n$ 时 $\delta\varepsilon/\varepsilon$ 还是不能成为大数.如果 $\eta n\gg1$，那么当 $\eta\sim1$ 时密度的相对变化达到 $C_{1}n^{2}$ 的量级，同时初始扰动的小量特性仅仅要求满足 $C_{1}n^{2}\eta_{0}^{2}\ll1$.这样一来，尽管扰动的增长进行得缓慢，但总的增加可能相当可观，结果扰动可能会比较大.

前面提到的第二和第三种类型的扰动可以采用类似的方式进行研究.然而从下述简单的理解出发，可以不通过计算细节就得到这些扰动的衰减律.

如果在不大的物质区域里（其线尺度为 $l$）有速度为 $\delta\boldsymbol{v}$ 的旋转扰动，那么这个区域的角动量 $\sim(\varepsilon/c^{2})l^{3}\cdot l\cdot v$.在宇宙膨胀时 $l$ 正比于 $a$ 增加，而 $\varepsilon$ 按照 $a^{-3}$（在 $p=0$ 的情况下）或者 $a^{-4}$（在 $p=\varepsilon/3$ 的情况下）减小.因此基于角动量守恒，我们有
\begin{equation}
  \delta v=\mathrm{const}\quad\text{当}\ p=\frac{\varepsilon}{3},
  \qquad
  \delta v\propto\frac{1}{a}\quad\text{当}\ p=0.
\end{equation}

最后，在宇宙膨胀时引力波的能量密度应该按照 $a^{-4}$ 减小.从另一方面，这个密度可通过度规的扰动表示为 $\sim k^{2}(h_{\alpha}^{\beta})^{2}$，其中 $k=n/a$ 为扰动的波矢.由此可得，引力波类型的扰动的振幅随时间按照 $1/a$ 的规律减小.

\section{均匀空间}

关于空间均匀性和各向同性的假设完全确定了它的度规（只剩下曲率的正负号可自由选取）.只假设均匀性这一个性质，而不带有任何附加的对称性，会留出大得多的自由度.我们将要考察的问题是，均匀空间的度规性质是怎样的.

我们将要讨论的是在给定时刻 $t$ 的空间度规.在假设时空参考系选取为同步的条件下，$t$ 对于整个空间就是统一的同步时间.

均匀性意味着在所有的空间点度规的性质都相同.这个概念的精确定义，涉及考虑一组能使空间变为自身，即保持其度规不变的坐标变换：如果在变换之前线元是
\begin{equation*}
  \mathrm{d}l^{2}=\gamma_{\alpha\beta}(x^{1},x^{2},x^{3})
  \mathrm{d}x^{\alpha}\mathrm{d}x^{\beta},
\end{equation*}
那么在变换之后同样的线元是
\begin{equation*}
  \mathrm{d}l^{2}=\gamma_{\alpha\beta}(x^{1},x^{2},x^{3})
  \mathrm{d}x^{\alpha}\mathrm{d}x^{\beta},
\end{equation*}
$\gamma_{\alpha\beta}$ 相对于新坐标具备同样的函数依赖关系.空间是均匀的，如果它允许存在一组变换（或称为运动群），使得它的任何一个给定点可以与其他任意的点相重合.基于空间的三维特性，显然，为了实现这一点，该群的不同变换应该用三个独立的参量来确定.

于是，在欧几里得空间里均匀性表现为度规相对于笛卡儿坐标系的平行移动（平移）的不变性.每个平移由三个参数确定——坐标原点的位移矢量的分量.所有这些变换保持构造线元的三个独立的微分 $(\mathrm{d}x,\mathrm{d}y,\mathrm{d}z)$ 不变.

在非欧几里得均匀空间的一般情况下，它的运动群的变换仍然使得三个独立的线性微分形式保持不变，但是无法简化为某一个坐标函数的全微分.我们将这些微分形式写成
\begin{equation}
  e_{\alpha}^{(a)}\,\mathrm{d}x^{\alpha},
\end{equation}
其中拉丁字母指标 $(a)$ 用于给三个独立的矢量（坐标的函数）编号，我们称这些矢量为一个标架.

借助于（116.1）式，相对于给定的运动群保持不变的空间度规可以这样建立
\begin{equation}
  \mathrm{d}l^{2}=\eta_{ab}
  (e_{\alpha}^{(a)}\mathrm{d}x^{\alpha})
  (e_{\beta}^{(b)}\mathrm{d}x^{\beta}),
\end{equation}
就是说，度规张量
\begin{equation}
  \gamma_{\alpha\beta}=\eta_{ab}e_{\alpha}^{(a)}e_{\beta}^{(b)},
\end{equation}
其中按指标 $a,b$ 对称的系数 $\eta_{ab}$ 是时间的函数.

这样一来，借助于三个标架矢量，我们得到了空间度规的“三维标架”表述；在§98节中得到的所有公式对这种表述都适用.标架矢量的选取取决于空间的对称性质，一般说来，这些矢量不是正交的（因此矩阵 $\eta_{ab}$ 不是对角的）.

如同在§98节中，与三个矢量 $e_{\alpha}^{(a)}$ 一起引入三个与它们互逆的矢量 $e_{(a)}^{\alpha}$.对于它们
\begin{equation}
  e_{(a)}^{\alpha}e_{\alpha}^{(b)}=\delta_{a}^{b},\qquad
  e_{(a)}^{\alpha}e_{\beta}^{(a)}=\delta_{\beta}^{\alpha}.
\end{equation}

在三维情况下，这两组矢量之间的联系可以表示为下列显性的形式，
\begin{equation}
  \boldsymbol{e}_{(1)}=\frac{1}{V}\boldsymbol{e}^{(2)}\times\boldsymbol{e}^{(3)},
  \quad
  \boldsymbol{e}_{(2)}=\frac{1}{V}\boldsymbol{e}^{(3)}\times\boldsymbol{e}^{(1)},
  \quad
  \boldsymbol{e}_{(3)}=\frac{1}{V}\boldsymbol{e}^{(1)}\times\boldsymbol{e}^{(2)},
\end{equation}
其中
\begin{equation*}
  V=|e_{\alpha}^{(a)}|
  =\boldsymbol{e}^{(1)}\cdot\boldsymbol{e}^{(2)}\times\boldsymbol{e}^{(3)},
\end{equation*}
而 $e_{(a)}$ 和 $e^{(a)}$ 应该理解为分别具有分量 $e_{(a)}^{\alpha}$ 和 $e_{\alpha}^{(a)}$ 的笛卡儿矢量.度规张量（116.3）的行列式
\begin{equation}
  \gamma=\eta v^{2},
\end{equation}
其中 $\eta$ 是矩阵 $\eta_{ab}$ 的行列式.

微分形式（116.1）的不变性意味着，
\begin{equation}
  e_{\alpha}^{(a)}(x)\mathrm{d}x^{\alpha}
  =e_{\alpha}^{(a)}(x)\mathrm{d}x^{\alpha},
\end{equation}
并且在等式两边的 $e_{\alpha}^{(a)}$ 分别是旧坐标和新坐标的同一个函数.给这个等式乘以 $e_{(a)}^{\beta}(x)$，替换 $\mathrm{d}x^{\beta}=(\partial x^{\beta}/\partial x^{\alpha})\mathrm{d}x^{\alpha}$ 并比较同一个微分 $\mathrm{d}x^{\alpha}$ 的系数，我们得到
\begin{equation}
  \frac{\partial x^{\beta}}{\partial x^{\alpha}}
  =e_{(a)}^{\beta}(x)e_{\alpha}^{(a)}(x).
\end{equation}

一系列微分方程组，决定了给定标架下的函数 $x^{\beta}(x)$\footnote{对于形式为 $x'^{\beta}=x^{\beta}+\xi^{\beta}$ 的变换，其中 $\xi^{\beta}$ 为小量，从（116.8）中可得到方程
\begin{equation}
\frac{\partial\xi^{\beta}}{\partial x^{\alpha}}
=\xi^{\gamma}e_{\alpha}^{(a)}\frac{\partial e_{(a)}^{\beta}}{\partial x^{\gamma}}.
\tag{116.8a}
\end{equation}
这些方程的三个线性独立解 $\xi_{(b)}^{\beta}$（$b=1,2,3$）确定空间运动群的无穷小变换. 矢量 $\xi_{(b)}^{\beta}$ 被称作基灵矢量（与 §94 第一个脚注相比较）.}.为了使之成为可积的，方程（116.8）应该恒定地满足条件
\begin{equation*}
  \frac{\partial^{2}x^{\beta}}{\partial x^{\alpha}\partial x^{\gamma}}
  =\frac{\partial^{2}x^{\beta}}{\partial x^{\gamma}\partial x^{\alpha}}.
\end{equation*}

算出导数，我们得到
\begin{equation*}
  \begin{aligned}
    &\left[\frac{\partial e_{(a)}^{\beta}(x)}{\partial x^{\delta}}
    e_{(b)}^{\delta}(x)
    -\frac{\partial e_{(b)}^{\beta}(x)}{\partial x^{\delta}}
    e_{(a)}^{\delta}(x)\right]
    e_{\gamma}^{(b)}(x)e_{\alpha}^{(a)}(x)=\\
    &\qquad=e_{(a)}^{\beta}(x)
    \left[\frac{\partial e_{\gamma}^{(a)}(x)}{\partial x^{\alpha}}
    -\frac{\partial e_{\alpha}^{(a)}(x)}{\partial x^{\gamma}}\right],
  \end{aligned}
\end{equation*}
给等式两边乘以 $e_{(d)}^{\alpha}(x)e_{(c)}^{\gamma}(x)e_{\beta}^{(f)}(x)$，使用（116.4）式将微分运算从一些因子上面转移到另一些因子上，在方程左边我们得到：
\begin{equation*}
  \begin{aligned}
    &e_{\beta}^{(f)}(x)
    \left[\frac{\partial e_{(d)}^{\beta}(x)}{\partial x^{\delta}}
    e_{(c)}^{\delta}(x)
    -\frac{\partial e_{(c)}^{\beta}(x)}{\partial x^{\delta}}
    e_{(d)}^{\delta}(x)\right]=\\
    &\qquad=e_{(c)}^{\beta}(x)e_{(d)}^{\delta}(x)
    \left[\frac{\partial e_{\beta}^{(f)}(x)}{\partial x^{\delta}}
    -\frac{\partial e_{\delta}^{(f)}(x)}{\partial x^{\beta}}\right],
  \end{aligned}
\end{equation*}
而在方程右边则得到关于 $x$ 的函数的同样的表达式.由于 $x$ 和 $x$ 是任意的，那么这些表达式应该化为常数：
\begin{equation}
  \left(\frac{\partial e_{\alpha}^{(c)}}{\partial x^{\beta}}
  -\frac{\partial e_{\beta}^{(c)}}{\partial x^{\alpha}}\right)
  e_{(a)}^{\alpha}e_{(b)}^{\beta}=C^{c}{}_{ab}.
\end{equation}

常数 $C^{c}{}_{ab}$ 被称做群的结构常数.两边同乘以 $e_{(c)}^{\gamma}$，可以将（116.9）重写为如下形式
\begin{equation}
  e_{(a)}^{\alpha}\frac{\partial e_{(b)}^{\gamma}}{\partial x^{\alpha}}
  -e_{(b)}^{\beta}\frac{\partial e_{(a)}^{\gamma}}{\partial x^{\beta}}
  =C^{c}{}_{ab}e_{(c)}^{\gamma}.
\end{equation}

这就是要寻找的空间均匀性的条件.等式（116.9）左边的表达式和量 $\lambda^{c}{}_{ab}$（98.10）的定义相同，这样一来，量 $\lambda^{c}{}_{ab}$ 是恒定的.

从结构常数的定义可以看出，它关于下标是反对称的：
\begin{equation}
  C^{c}{}_{ab}=-C^{c}{}_{ba}.
\end{equation}

对于它们还能得到一个条件，首先指出，等式（116.10）可以写成对易关系的形式：
\begin{equation}
  [X_{a},X_{b}]\equiv X_{a}X_{b}-X_{b}X_{a}=C^{c}{}_{ab}X_{c}
\end{equation}
其中线性微分算子\footnote{在连续群（或称李群）的数学理论中，满足（116.12）形式的条件的算子称作群的生成元. 为了避免在与其他表示相比较时产生误解，我们还是要指出，连续群的系统理论通常是以基灵矢量确定的算子 $X_{a}=\xi_{(a)}^{\alpha}\partial/\partial x^{\alpha}$ 为出发点而建立的.}
\begin{equation}
  X_{a}=e_{(a)}^{\alpha}\frac{\partial}{\partial x^{\alpha}}.
\end{equation}

于是上面提到的关系可从下列恒等式中产生
\begin{equation*}
  [[X_{a},X_{b}],X_{c}]+[[X_{b},X_{c}],X_{a}]
  +[[X_{c},X_{a}],X_{b}]=0
\end{equation*}
（这就是雅可比恒等式），并且具有形式
\begin{equation}
  C^{f}{}_{ab}C^{d}{}_{cf}
  +C^{f}{}_{bc}C^{d}{}_{af}
  +C^{f}{}_{ca}C^{d}{}_{bf}=0.
\end{equation}

使用一组双指标量来代替三指标常数 $C^{c}{}_{ab}$ 会呈现出一定的优点.这些双指标量通过如下对偶变换得到：
\begin{equation}
  C^{c}{}_{ab}=e_{abd}C^{dc},
\end{equation}
其中 $e_{abc}=e^{abc}$ 为单位反对称符号（其中 $e_{123}=+1$）.利用这些常数，对易关系（116.12）可写为如下形式
\begin{equation}
  e^{abc}X_{b}X_{c}=C^{ad}X_{d}.
\end{equation}

性质（116.16）在定义（116.15）中已经考虑到了，而性质（116.14）可以采取如下形式
\begin{equation}
  e_{bcd}C^{cd}C^{ba}=0.
\end{equation}

同时我们指出，对于量 $C^{ab}$ 定义（116.9）可以表示为矢量形式
\begin{equation}
  C^{ab}=-\frac{1}{v}\boldsymbol{e}^{(a)}\cdot
  \operatorname{rot}\boldsymbol{e}^{(b)},
\end{equation}
其中进行矢量运算时，仍旧把坐标 $x^{\alpha}$ 视同为笛卡儿坐标.

在微分形式（116.1）（连同它们一起的还有算子 $X_{a}$）中三个标架矢量的选取当然不是唯一的.可以对它们施以任意的、带有恒定系数的线性变换：
\begin{equation}
  \boldsymbol{e}_{(a)}=A_{a}^{b}\boldsymbol{e}_{(b)}.
\end{equation}

在这些变换下，量 $\eta_{ab}$ 和 $C^{ab}$ 表现出类似于张量的性质.

条件（116.17）是结构常数 $C^{ab}$ 应该满足的唯一条件.但是在满足这些条件的常数中存在着若干等价组，同一组中的常数之间的区别仅仅和变换（116.19）有关.均匀空间的分类问题可归结为确定所有的不等价的结构常数组.利用量 $C^{ab}$ 的“张量”性质，通过下面的简单的方法可以做到这一点（C. G. Behr, 1962）.

不对称的“张量”$C^{ab}$ 可以分解为对称的和反对称的部分.第一部分表示为 $n^{ab}$，而第二部分通过它的对偶“矢量”$a_{c}$ 表示：
\begin{equation}
  C^{ab}=n^{ab}+e^{abc}a_{c}.
\end{equation}

这个表达式代入（116.17）会导出条件
\begin{equation}
  n^{ab}a_{b}=0.
\end{equation}

利用变换（116.19）可以将对称“张量”$n^{ab}$ 化为对角形式；设 $n_{1},n_{2},n_{3}$ 是它的主值.等式（116.21）表明，“矢量”$a_{b}$（如果它存在）位于“张量”$n^{ab}$ 的主方向之一，即对应于零主值的那个方向.所以，不失一般性，可以令 $a_{b}=(a,0,0)$，那么（116.21）可化为 $an_{1}=0$，就是说量 $a$ 或者 $n_{1}$ 的其中之一应该为零.对易关系（116.16）可取如下形式
\begin{equation}
  \begin{gathered}
    [X_{1},X_{2}]=-aX_{2}+n_{3}X_{3},\\
    [X_{2},X_{3}]=n_{1}X_{1},\\
    [X_{3},X_{1}]=n_{2}X_{2}+aX_{3}.
  \end{gathered}
\end{equation}

在这之后还剩下的自由是改变算子 $X_{a}$ 的正负号以及它们的任意标度变换（乘以常数）.这种自由允许我们同时改变所有的 $n_{1},n_{2},n_{3}$ 的符号，同时使得量 $a$ 的符号为正（如果它不为零）.我们还可以将所有的结构常数化为 $\pm1$，只要量 $a,n_{2},n_{3}$ 其中之一为零.如果这三个量全都不为零，那么在标度变换下商 $a^{2}/n_{2}n_{3}$ 为不变量\footnote{严格说来，为了使 $C^{ab}$ 的张量特性得到遵循，应该在定义（116.15）中引入因子 $\sqrt{\eta}$（比较 §83 节中提到的关于应该如何相对于任意的坐标变换来确定反对称的单位张量）. 但是这里我们将不再深入到这些细节中：为了我们的目的，我们可以直接从（116.22）中提取出结构常数的变换规律.}.

这样，我们得到均匀空间可能类型的列表如下；表格第一列中的罗马数字表示按照比安基分类法而得的类型的编号（L. Bianchi, 1918）：\footnote{参数 $a$ 可以取任何正数值. 相应的类型事实上是不同群的单参数族，它们被赋予类型 VI 和 VII 只是出于惯例.}.

\begin{center}
\begin{tabular}{ccccc}
\hline
 类型 & $a$ & $n_{1}$ & $n_{2}$ & $n_{3}$ \\
\hline
 I & 0 & 0 & 0 & 0 \\
 II & 0 & 1 & 0 & 0 \\
 $\mathrm{VII}_{0}$ & 0 & 1 & 1 & 0 \\
 $\mathrm{VI}_{0}$ & 0 & 1 & $-1$ & 0 \\
 IX & 0 & 1 & 1 & 1 \\
 VIII & 0 & 1 & 1 & $-1$ \\
 V & 1 & 0 & 0 & 0 \\
 IV & 1 & 0 & 0 & 1 \\
 $\mathrm{VII}_{a}$ & $a$ & 0 & 1 & 1 \\
 $\left.\begin{array}{l}
   \mathrm{III}\ (a=1)\\
   \mathrm{VI}_{a}\ (a\neq1)
 \end{array}\right\}$ & $a$ & 0 & 1 & $-1$ \\
\hline
\end{tabular}
\end{center}

类型 I 是欧几里得空间；空间曲率张量（参考下文的公式（116.24））的所有分量为零.除了伽利略度规的平凡情形，这里还涉及含时度规，后者将在下一节进行研究.

作为一个特例，恒定正曲率的空间包含于类型 IX 之中.如果在线元（116.2）中设定 $\eta_{ab}=\delta_{ab}/4\lambda$ 就可以得到它，其中 $\lambda$ 为正的常数.实际上，按照（116.24）和条件 $C^{11}=C^{22}=C^{33}=1$（类型 IX 的结构常数）的计算给出 $P_{(a)(b)}=(1/2)\delta_{ab}$，然后有
\begin{equation*}
  P_{\alpha\beta}=P_{(a)(b)}e_{\alpha}^{(a)}e_{\beta}^{(b)}
  =2\lambda\gamma_{\alpha\beta},
\end{equation*}
这正好对应于上面指出的空间（比较（111.3））.

恒定负曲率空间以相似的方式作为一个特例包含在类型 V 中.实际上，设定 $\eta_{ab}=\delta_{ab}/\lambda$ 并且按照（116.24）和条件 $C^{23}=-C^{32}=1$ 算出 $P_{(a)(b)}$，我们得到
\begin{equation*}
  P_{(a)(b)}=-2\delta_{ab},\qquad
  P_{\alpha\beta}=-2\lambda\gamma_{\alpha\beta},
\end{equation*}
这对应于恒定负曲率.

最后，我们展示一下均匀空间宇宙的爱因斯坦方程是如何归结为只包含时间函数的常微分方程组的.为了这一目的，我们必须将四维矢量和四维张量沿着空间三维标架的基矢方向分解出空间分量：
\begin{equation*}
  R_{(a)(b)}=R_{\alpha\beta}e_{(a)}^{\alpha}e_{(b)}^{\beta},\qquad
  R_{0(a)}=R_{0\alpha}e_{(a)}^{\alpha},\qquad
  u^{(a)}=u^{\alpha}e_{\alpha}^{(a)},
\end{equation*}
所有的这些量现在都只是时间 $t$ 的函数；能量密度 $\varepsilon$ 和物质的压强 $p$ 这样的标量同样也是时间的函数.

根据（97.11）—（97.13），同步参考系中的爱因斯坦方程可通过三维张量 $\varkappa_{\alpha\beta}$ 和 $P_{\alpha\beta}$ 来表达.第一步我们只有
\begin{equation}
  \varkappa_{(a)(b)}=\dot{\eta}_{ab},\qquad
  \varkappa_{(a)}^{(b)}=\dot{\eta}_{ac}\eta^{cb}
\end{equation}
（符号上方的点表示对 $t$ 微分）.$P_{(a)(b)}$ 的分量可以利用（98.14）通过量 $\eta_{ab}$ 和群的结构常数来表达.在将三指标的 $\lambda^{a}{}_{bc}=C^{a}{}_{bc}$ 替换为双指标的 $C^{ab}$ 并且经过一系列变换\footnote{其中使用了公式 $\eta_{ad}\eta_{be}\eta_{cf}e^{def}=\eta e_{abc}$，$e_{abf}e^{cdf}=\delta_{a}^{c}\delta_{b}^{d}-\delta_{a}^{d}\delta_{b}^{c}$.}之后，我们得到
\begin{equation}
  P_{(a)}^{(b)}=\frac{1}{2\eta}\bigl\{
  2C^{bd}C_{ad}+C^{db}C_{ad}+C^{bd}C_{da}
  -C^{d}{}_{d}(C^{b}{}_{a}+C_{a}{}^{b})
  +\delta_{a}^{b}[(C^{d}{}_{d})^{2}-2C^{df}C_{df}]\bigr\}.
\end{equation}

这里，按照一般的规则
\begin{equation*}
  C_{a}{}^{b}=\eta_{ac}C^{cb},\qquad
  C_{ab}=\eta_{ac}\eta_{bd}C^{cd}.
\end{equation*}

同时指出，在均匀空间中三维张量 $P_{\alpha\beta}$ 的比安基恒等式具有如下形式
\begin{equation}
  P_{b}^{c}C^{b}{}_{ca}+P_{a}^{c}C^{b}{}_{cb}=0.
\end{equation}

四维里奇张量\footnote{包含在 $R_{\alpha}^{0}$ 中的协变导数 $\varkappa_{\alpha;\gamma}^{\beta}$ 可利用 §98 最后一个脚注中列出的公式进行换算.}的三维标架分量的最终表达式如下：
\begin{equation}
  \begin{gathered}
    R_{0}^{0}=-\frac{1}{2}\dot{\varkappa}_{(a)}^{(a)}
    -\frac{1}{4}\varkappa_{(a)}^{(b)}\varkappa_{(b)}^{(a)},\\
    R_{(a)}^{0}
    =-\frac{1}{2}\varkappa_{(b)}^{(c)}
    (C^{b}{}_{ca}-\delta_{a}^{b}C^{d}{}_{dc}),\\
    R_{(a)}^{(b)}
    =-\frac{1}{2\sqrt{\eta}}
    \dot{(\sqrt{\eta}\varkappa_{(a)}^{(b)})}
    -P_{(a)}^{(b)}.
  \end{gathered}
\end{equation}

我们强调，这样一来，在建构爱因斯坦方程时，没有必要使用标架矢量作为坐标函数那样的显式表达式.

\section{平直各向异性模型}

各向同性模型对描述宇宙后期演化是适用的，但不能据此预期，它也适用于描述宇宙靠近时间奇点时的早期演化.这个问题将在§119中详细讨论，而在本节和下一节将对爱因斯坦方程的一些解进行初步研究，这些解同样具有时间奇点，但原则上属于不同的类型（有别于弗里德曼奇异性）.

我们将寻找这样的解，其中度规张量的所有分量在适当选取的参考系下仅仅是一个变量——时间 $x^{0}=t$ 的函数\footnote{在 §117，§118 节中为了简化公式的书写我们设定 $c=1$.}.这样的问题已经在§109中研究过，但是，那里只研究了当行列式 $|g_{\alpha\beta}|=0$ 时的情况.现在我们将认为这个行列式不为零.就像在§109中讲解的那样，在这种情况下可以不失一般性地令所有的 $g_{0\alpha}=0$.通过对变量 $t$ 做 $\sqrt{g_{00}}\,\mathrm{d}t\to\mathrm{d}t$ 的变换，可以使 $g_{00}$ 等于1，于是我们得到同步参考系，在其中
\begin{equation}
  g_{00}=1,\quad g_{0\alpha}=0,\quad
  g_{\alpha\beta}=-\gamma_{\alpha\beta}(t).
\end{equation}

现在我们可以利用形式为（97.11）—（97.13）的爱因斯坦方程.由于量 $\gamma_{\alpha\beta}$，以及随之三维张量 $\varkappa_{\alpha\beta}=\dot{\gamma}_{\alpha\beta}$ 的分量不依赖于坐标 $x^{\alpha}$，于是 $R_{0\alpha}\equiv0$.按照同样的原因 $P_{\alpha\beta}\equiv0$，结果真空中的引力场方程可化为下列方程组：
\begin{equation}
  \dot{\varkappa}_{\alpha}^{\alpha}
  +\frac{1}{2}\varkappa_{\alpha}^{\beta}\varkappa_{\beta}^{\alpha}=0,
\end{equation}
\begin{equation}
  \frac{1}{\sqrt{\gamma}}
  (\sqrt{\gamma}\varkappa_{\alpha}^{\beta})^{\cdot}=0.
\end{equation}

从（117.3）可导出
\begin{equation}
  \sqrt{\gamma}\varkappa_{\alpha}^{\beta}=2\lambda_{\alpha}^{\beta},
\end{equation}
其中 $\lambda_{\alpha}^{\beta}$ 为恒定值.对指标 $\alpha$ 和 $\beta$ 进行缩并，我们得到
\begin{equation*}
  \varkappa_{\alpha}^{\alpha}=\frac{\dot{\gamma}}{\gamma}
  =\frac{2}{\sqrt{\gamma}}\lambda_{\alpha}^{\alpha},
\end{equation*}
由此可见，$\gamma=\mathrm{const}\cdot t^{2}$.不失一般性可以令 $\mathrm{const}=1$（这一点可以简单地通过改变坐标 $x^{\alpha}$ 的标度而达到）；那么 $\lambda_{\alpha}^{\alpha}=1$.现在将（117.4）代入方程（117.2）可给出关系式
\begin{equation}
  \lambda_{\alpha}^{\beta}\lambda_{\beta}^{\alpha}=1,
\end{equation}
该式将常数 $\lambda_{\alpha}^{\beta}$ 彼此之间关联起来.

接下来，将（117.4）中的指标 $\beta$ 下降，将这些等式写成关于 $\gamma_{\alpha\beta}$ 的常微分方程组的形式：
\begin{equation}
  \dot{\gamma}_{\alpha\beta}
  =\frac{2}{t}\lambda_{\alpha}^{\gamma}\gamma_{\gamma\beta}.
\end{equation}

系数 $\lambda_{\alpha}^{\gamma}$ 的集合可以看做某种线性代换的矩阵.通过坐标 $x^{1},x^{2},x^{3}$（或者等效地，量 $g_{1\beta},g_{2\beta},g_{3\beta}$）的相应线性变换，一般说来可以将这个矩阵化为对角形式.我们将它的主值表示为 $p_{1},p_{2},p_{3}$，并且认为它们全都是实数并且是不同的（其他情况参考下文）；相应的主轴上的单位矢量是 $\boldsymbol{n}^{(1)},\boldsymbol{n}^{(2)},\boldsymbol{n}^{(3)}$.那么方程（117.6）的解可以表示为如下形式
\begin{equation}
  \gamma_{\alpha\beta}
  =t^{2p_{1}}n_{\alpha}^{(1)}n_{\beta}^{(1)}
  +t^{2p_{2}}n_{\alpha}^{(2)}n_{\beta}^{(2)}
  +t^{2p_{3}}n_{\alpha}^{(3)}n_{\beta}^{(3)}
\end{equation}
（通过对坐标作合适的标度变换，可将 $t$ 的幂的常系数化为1）.最后选取矢量 $\boldsymbol{n}^{(1)},\boldsymbol{n}^{(2)},\boldsymbol{n}^{(3)}$ 的方向作为坐标轴（将其称为 $x,y,z$）的方向，我们最终将度规化为
\begin{equation}
  \mathrm{d}s^{2}=\mathrm{d}t^{2}
  -t^{2p_{1}}\mathrm{d}x^{2}
  -t^{2p_{2}}\mathrm{d}y^{2}
  -t^{2p_{3}}\mathrm{d}z^{2}
\end{equation}
（E. Kasner, 1922）.这里 $p_{1},p_{2},p_{3}$ 是满足下面两个关系式的任意三个数：
\begin{equation}
  p_{1}+p_{2}+p_{3}=1,\qquad
  p_{1}^{2}+p_{2}^{2}+p_{3}^{2}=1
\end{equation}
（第一个由 $-g=t^{2}$ 导出，而第二个由（117.5）得到）.

显然，$p_{1},p_{2},p_{3}$ 三个数不可能都有相同的值.他们当中有两个相等的情况是：取值为 $(0,0,1)$ 或 $(-1/3,2/3,2/3)$.在其他任何情况下 $p_{1},p_{2},p_{3}$ 是不同的，并且其中的一个为负，而另外两个为正.如果以 $p_{1}<p_{2}<p_{3}$ 的顺序排列，那么它们的值将分别位于区间：
\begin{equation}
  -\frac{1}{3}\leqslant p_{1}\leqslant0,\qquad
  0\leqslant p_{2}\leqslant\frac{2}{3},\qquad
  \frac{2}{3}\leqslant p_{3}\leqslant1.
\end{equation}

这样一来，度规（117.8）对应于平直、均匀、但各向异性的空间，其中所有的体积都（随着时间的增加）正比于 $t$ 增长，并且沿两个轴 $(y,z)$ 的线性距离增加，而沿一个轴 $(x)$ 的线性距离减小.时刻 $t=0$ 是解的奇点；在该点度规有奇异性，这个奇异性通过任何参考系变换都无法消除，并且四维曲率张量的不变量趋于无穷大.只有 $p_{1}=p_{2}=0,\ p_{3}=1$ 的情况是个例外，这时我们只不过是在和平直时空打交道：度规（117.8）经过变换 $t\sinh z=\zeta,\ t\cosh z=\tau$ 化为伽利略度规\footnote{（117.8）的参数为类空的情况下也存在同样类型的解；此时我们只需适当地改变其中的符号，例如：
\[
\mathrm{d}s^{2}=x^{2p_{1}}\mathrm{d}t^{2}-\mathrm{d}x^{2}-x^{2p_{2}}\mathrm{d}y^{2}-x^{2p_{3}}\mathrm{d}z^{2}.
\]
但是在这种情况下，同时还存在另一种类型的解，这时方程（117.6）中的矩阵 $\lambda_{\alpha}^{\beta}$ 有复数主值或相同的主值（参考习题 1 和 2）. 在类时参数 $t$ 的情况下这些解是不可能存在的，因为其中的行列式 $g$ 不满足必要条件 $g<0$.
我们同时给出相关的参考文献，其中求得了真空中爱因斯坦方程各种不同类型的严格解，它们与更多的参量相关：B. K. Harrison//Phys. Rev. 1959. V. 116. P. 1285.}.

度规（117.8）是真空中爱因斯坦方程的严格解.而在奇点附近，在 $t$ 很小的时候，它也是物质在空间中均匀分布情况的近似解（精度到 $1/t$ 阶）.物质密度变化的速度和过程在这种情况下只取决于物质在给定引力场中的运动方程，而物质对场的反作用可以忽略不计.在 $t\to0$ 时物质的密度趋向无穷大——对应于奇异性的物理特征（参考习题3）.

\begin{xiti}

\noindent{\heiti 习题1}\quad 在矩阵 $\lambda_{\alpha}^{\beta}$ 具有一个实数（$p_{3}$）和两个复数（$p_{1,2}=p\pm\mathrm{i}p$）主值的情况下，求方程（117.6）的解.

解：在这种情况下，所有的量都依赖于变量 $x^{0}$，后者应该具有类空的特征；我们将其表示为 $x^{0}=x$.所以，在（117.1）中现在应该有 $g_{00}=-1$.方程（117.2）和（117.3）无变化.

在（117.7）中矢量 $\boldsymbol{n}^{(1)},\boldsymbol{n}^{(2)}$ 成为复矢量：$\boldsymbol{n}^{(1,2)}=(\boldsymbol{n}\pm\mathrm{i}\boldsymbol{n})/\sqrt{2}$，其中 $\boldsymbol{n},\boldsymbol{n}$ 为单位矢量.沿 $\boldsymbol{n},\boldsymbol{n},\boldsymbol{n}^{(3)}$ 的方向选取坐标轴 $x^{1},x^{2},x^{3}$，我们得到下列形式的解
\begin{equation*}
  \begin{gathered}
    -g_{11}=g_{22}=x^{2p}\cos\left(2p\ln\frac{x}{a}\right),\qquad
    g_{12}=-x^{2p}\sin\left(2p\ln\frac{x}{a}\right),\\
    g_{33}=-x^{2p_{3}},\qquad
    -g=-g_{00}|g_{\alpha\beta}|=x^{2},
  \end{gathered}
\end{equation*}
其中 $a$ 为常数（无法在只对 $x$ 轴作标度变换、同时又不改变上述表达式中其他系数的情况下消除 $a$）.$p_{1},p_{2},p_{3}$ 三个数照旧满足关系式（117.9），并且实数 $p_{3}$ 要么小于 $-1/3$，要么大于1.

\noindent{\heiti 习题2}\quad 在两个主值相同（$p_{2}=p_{3}$）的情况下求解同样的问题.

解：我们知道，根据线性微分方程的一般理论，在这种情况下方程组（117.6）可以化为下列正则形式：
\begin{equation*}
  \dot{g}_{11}=\frac{2p_{1}}{x}g_{11},\qquad
  \dot{g}_{2\alpha}=\frac{2p_{2}}{x}g_{2\alpha},\qquad
  \dot{g}_{3\alpha}=\frac{2p_{2}}{x}g_{3\alpha}
  +\frac{\lambda}{x}g_{2\alpha},\qquad\alpha=2,3,
\end{equation*}
其中 $\lambda$ 是常数.在 $\lambda=0$ 时我们回到（117.8）.在 $\lambda\neq0$ 时可以设定 $\lambda=1$；那么
\begin{equation*}
  g_{11}=-x^{2p_{1}},\qquad
  g_{2\alpha}=a_{\alpha}x^{2p_{2}},\qquad
  g_{3\alpha}=a_{\alpha}x^{2p_{2}}\ln x+b_{\alpha}x^{2p_{2}}.
\end{equation*}
从条件 $g_{32}=g_{23}$ 求得 $a_{2}=0,\ a_{3}=b_{2}$.通过适当选取坐标轴 $x^{2},x^{3}$ 的标度我们最终将度规化为下列形式：
\begin{equation*}
  \mathrm{d}s^{2}=-\mathrm{d}x^{2}
  -x^{2p_{1}}(\mathrm{d}x^{1})^{2}
  \pm2x^{2p_{2}}\mathrm{d}x^{2}\mathrm{d}x^{3}
  \pm x^{2p_{2}}\ln\frac{x}{a}(\mathrm{d}x^{3})^{2}.
\end{equation*}
$p_{1},p_{2}$ 可以具有数值 $1,0$ 或者 $-1/3,2/3$.

\noindent{\heiti 习题3}\quad 设物质均匀分布在度规为（117.8）的空间内，求物质的密度在奇点 $t=0$ 附近随时间的变化规律.

解：忽略物质对场的反作用，从流体力学的运动方程出发
\begin{equation*}
  \frac{1}{\sqrt{-g}}\frac{\partial}{\partial x^{i}}
  (\sqrt{-g}\,\sigma u^{i})=0,
\end{equation*}
\begin{equation}\tag{1}
  (p+\varepsilon)u^{k}
  \left(\frac{\partial u_{i}}{\partial x^{k}}
  -\frac{1}{2}u^{l}\frac{\partial g_{kl}}{\partial x^{i}}\right)
  =-\frac{\partial p}{\partial x^{i}}
  -u_{i}u^{k}\frac{\partial p}{\partial x^{k}},
\end{equation}
这些方程包含在方程 $T_{i;k}^{k}=0$ 中（参考本教程第六卷§134）.这里 $\sigma$ 是熵密度；在奇点附近应该使用极端相对论物态方程 $p=\varepsilon/3$，那么 $\sigma\propto\varepsilon^{3/4}$.

我们将（117.8）中的时间因子表示为 $a=t^{p_{1}},\ b=t^{p_{2}},\ c=t^{p_{3}}$.由于所有的量只依赖于时间，$\sqrt{-g}=abc$，方程（1）给出
\begin{equation*}
  \frac{\mathrm{d}}{\mathrm{d}t}(abc\,u_{0}\varepsilon^{3/4})=0,
  \qquad
  4\varepsilon\frac{\mathrm{d}u_{\alpha}}{\mathrm{d}t}
  +u_{\alpha}\frac{\mathrm{d}\varepsilon}{\mathrm{d}t}=0.
\end{equation*}
由此
\begin{equation}\tag{2}
  abc\,u_{0}\varepsilon^{3/4}=\mathrm{const},
\end{equation}
\begin{equation}\tag{3}
  u_{\alpha}\varepsilon^{1/4}=\mathrm{const}.
\end{equation}

按照（3），所有的协变分量 $u_{\alpha}$ 为相同阶数的量.在逆变分量中最大的（当 $t\to0$ 时）是 $u^{3}\approx u_{3}/c^{2}$.在恒等式 $u_{i}u^{i}=1$ 中仅保留最大的那些项，从而得到 $u_{0}^{2}\approx u_{3}u^{3}=(u_{3})^{2}/c^{2}$，然后由（2）和（3）可得：
\begin{equation*}
  \varepsilon\sim\frac{1}{a^{2}b^{2}},\qquad
  u_{\alpha}\sim\sqrt{ab},
\end{equation*}
或者
\begin{equation}\tag{4}
  \varepsilon\sim t^{-2(p_{1}+p_{2})}=t^{-2(1-p_{3})},\qquad
  u_{\alpha}\sim t^{(1-p_{3})/2}.
\end{equation}

这样可以得出，对于 $p_{3}$ 的所有的值，当 $t\to0$ 时 $\varepsilon$ 趋向无穷大，只有当 $p_{3}=1$ 是个例外，——与之对应的事实是，带有指数 $(0,0,1)$ 的度规的奇异性是虚假的.

所采用的近似方法的正确性可通过估算分量 $T_{i}^{k}$ 来检验，后者在方程（117.2），（117.3）的右边被忽略掉了.它们当中的主项为：
\begin{equation*}
  \begin{aligned}
    &T_{0}^{0}\sim\varepsilon u_{0}^{2}\sim t^{-(1+p_{3})},\qquad
    T_{1}^{1}\sim\varepsilon\sim t^{-2(1-p_{3})},\\
    &T_{2}^{2}\sim\varepsilon u_{2}u^{2}\sim t^{-(1+2p_{2}-p_{3})},\qquad
    T_{3}^{3}\sim\varepsilon u_{3}u^{3}\sim t^{-(1+p_{3})}.
  \end{aligned}
\end{equation*}
实际上在 $t\to0$ 时它们全都增长得比方程左边慢，后者以 $t^{-2}$ 增长.

\end{xiti}

\section{靠近奇点的振动状态}

我们将以第IX类均匀空间宇宙模型为例，来研究具有振动特性的度规的时间奇异性（В. А. Белинский, Е. М. Лифшиц, И. М. Халатников, 1968）.在下一节中我们将看到，这种特性具有相当普遍的意义.

我们感兴趣的是模型在奇点（我们选作时间原点 $t=0$）附近的表现.正如在§117节研究的卡兹涅尔（Kasner）解中所反映的那样，物质的存在不会对模型表现出的本质属性造成影响，为了简化研究过程，我们将假设空间是空的.

在（116.3）中将量 $\eta_{ab}(t)$ 的矩阵设定为对角的，将其对角元表示为 $a^{2},b^{2},c^{2}$；三个标架矢量 $\boldsymbol{e}^{(1)},\boldsymbol{e}^{(2)},\boldsymbol{e}^{(3)}$ 现在表示为 $\boldsymbol{l},\boldsymbol{m},\boldsymbol{n}$.那么空间度规可写成如下形式
\begin{equation}
  \gamma_{\alpha\beta}
  =a^{2}l_{\alpha}l_{\beta}
  +b^{2}m_{\alpha}m_{\beta}
  +c^{2}n_{\alpha}n_{\beta}.
\end{equation}

对于类型IX的空间，结构常数如下\footnote{对应于这些常数的标架矢量：
\[
\boldsymbol{l}=(\sin x^{3},-\cos x^{3}\sin x^{1},0),\qquad
\boldsymbol{m}=(\cos x^{3},\sin x^{3}\sin x^{1},0),\qquad
\boldsymbol{n}=(0,\cos x^{1},1).
\]
坐标的取值范围在区间 $0\leqslant x^{1}\leqslant\pi$，$0\leqslant x^{2}\leqslant2\pi$，$0\leqslant x^{3}\leqslant4\pi$ 内. 空间是闭合的，并且它的体积
\[
V=\int\sqrt{\gamma}\,\mathrm{d}x^{1}\mathrm{d}x^{2}\mathrm{d}x^{3}
=abc\int\sin x^{1}\,\mathrm{d}x^{1}\mathrm{d}x^{2}\mathrm{d}x^{3}
=16\pi^{2}abc.
\]
当 $a=b=c$ 时它会过渡到曲率半径为 $2a$ 的恒定正曲率空间.}：
\begin{equation}
  C^{11}=C^{22}=C^{33}=1
\end{equation}
（和 $C^{1}{}_{23}=C^{2}{}_{31}=C^{3}{}_{12}=1$）.

从（116.26）可以看出，对于这些常数，以及对角矩阵 $\eta_{ab}$，里奇张量的分量 $R_{(a)}^{0}$ 在同步参考系中恒为零.按照（116.24）非对角分量 $P_{(a)(b)}$ 同时也为零.爱因斯坦方程的剩余成分给出针对函数 $a(t),b(t),c(t)$ 的下列方程组：
\begin{equation}
  \begin{gathered}
    \frac{(\dot{a}bc)^{\cdot}}{abc}
    =\frac{1}{2a^{2}b^{2}c^{2}}[(b^{2}-c^{2})^{2}-a^{4}],\\
    \frac{(a\dot{b}c)^{\cdot}}{abc}
    =\frac{1}{2a^{2}b^{2}c^{2}}[(a^{2}-c^{2})^{2}-b^{4}],\\
    \frac{(ab\dot{c})^{\cdot}}{abc}
    =\frac{1}{2a^{2}b^{2}c^{2}}[(a^{2}-b^{2})^{2}-c^{4}],
  \end{gathered}
\end{equation}
\begin{equation}
  \frac{\ddot{a}}{a}+\frac{\ddot{b}}{b}+\frac{\ddot{c}}{c}=0
\end{equation}
（（118.3）是方程组 $R_{(1)}^{(1)}=R_{(2)}^{(2)}=R_{(3)}^{(3)}=0$；（118.4）是方程 $R_{(0)}^{(0)}=0$）.

在方程组（118.3），（118.4）中，对时间的导数项可以有更简单的形式，只要引入函数 $a,b,c$ 的对数 $\alpha,\beta,\gamma$ 来代替它们：
\begin{equation}
  a=\mathrm{e}^{\alpha},\quad b=\mathrm{e}^{\beta},\quad c=\mathrm{e}^{\gamma},
\end{equation}
而按下式用变量 $\tau$ 代替 $t$
\begin{equation}
  \mathrm{d}t=abc\,\mathrm{d}\tau.
\end{equation}
那么：
\begin{equation}
  \begin{gathered}
    2\alpha_{,\tau,\tau}=(b^{2}-c^{2})^{2}-a^{4},\\
    2\beta_{,\tau,\tau}=(a^{2}-c^{2})^{2}-b^{4},\\
    2\gamma_{,\tau,\tau}=(a^{2}-b^{2})^{2}-c^{4},
  \end{gathered}
\end{equation}
\begin{equation}
  \frac{1}{2}(\alpha+\beta+\gamma)_{,\tau,\tau}
  =\alpha_{,\tau}\beta_{,\tau}
  +\alpha_{,\tau}\gamma_{,\tau}
  +\beta_{,\tau}\gamma_{,\tau},
\end{equation}
其中下标 $,\tau$ 表示对 $\tau$ 微分.将（118.7）的各方程相加，并用（118.8）替换左边的二阶导数之和，我们得到
\begin{equation}
  \alpha_{,\tau}\beta_{,\tau}
  +\alpha_{,\tau}\gamma_{,\tau}
  +\beta_{,\tau}\gamma_{,\tau}
  =\frac{1}{4}(a^{4}+b^{4}+c^{4}
  -2a^{2}b^{2}-2a^{2}c^{2}-2b^{2}c^{2}).
\end{equation}

这个关系式只包含一阶导数并且是方程（118.7）的初积分.

方程（118.3）和（118.4）无法用解析形式严格求解，但在奇点附近可以进行详细的定性研究.

首先我们指出，当方程（118.3）或（118.7）的右边不存在时，方程组具有严格解，在其中
\begin{equation}
  a\sim t^{p_{l}},\qquad b\sim t^{p_{m}},\qquad c\sim t^{p_{n}},
\end{equation}
其中 $p_{l},p_{m},p_{n}$ 是通过下面的关系式联系起来的数：
\begin{equation}
  p_{l}+p_{m}+p_{n}
  =p_{l}^{2}+p_{m}^{2}+p_{n}^{2}=1
\end{equation}
（类似于均匀平直空间的卡兹涅尔解（117.8））.这里我们将幂指数表示为 $p_{l},p_{m},p_{n}$，并没有确定他们的大小顺序；我们再次使用§117中 $p_{1},p_{2},p_{3}$ 的表示法，将这三个数按照 $p_{1}<p_{2}<p_{3}$ 的大小顺序排列，它们的取值范围见（117.10）.这些数可以表示为下列参数形式
\begin{equation}
  p_{1}(u)=\frac{-u}{1+u+u^{2}},\qquad
  p_{2}(u)=\frac{1+u}{1+u+u^{2}},\qquad
  p_{3}(u)=\frac{u(1+u)}{1+u+u^{2}}.
\end{equation}

如果参数 $u$ 在 $u\geqslant1$ 的全范围内取值，则可以得到 $p_{1},p_{2},p_{3}$ 的所有不同的值（保持大小顺序不变）.$u<1$ 时同样的数值范围可通过下列替换得到：
\begin{equation}
  p_{1}\left(\frac{1}{u}\right)=p_{1}(u),\qquad
  p_{2}\left(\frac{1}{u}\right)=p_{3}(u),\qquad
  p_{3}\left(\frac{1}{u}\right)=p_{2}(u).
\end{equation}

在图25中描绘了 $p_{1},p_{2},p_{3}$ 依赖于 $1/u$ 的曲线.

\par\nopagebreak\noindent\begin{minipage}{\linewidth}\centering\includegraphics[width=6cm]{fig25.pdf}\\[0.3em]{\small 图 25}\end{minipage}\par\nopagebreak

我们假设在某个时段方程（118.7）的右边确实很小，使得它们可以被忽略，因而有卡兹涅尔状态（118.10）.但是当 $t\to0$ 时这样的情况不能无限制地持续，因为在已指出的各项当中总是存在增长的项.如果负的幂指数对应函数 $a(t)$（$p_{l}=p_{1}<0$），那么卡兹涅尔状态的扰动从 $a^{4}$ 项中产生，剩余的项在 $t$ 减小时会下降.

在方程（118.7）的右边仅保留这些扰动项，我们得到方程组：
\begin{equation}
  \alpha_{,\tau,\tau}=-\frac{1}{2}\mathrm{e}^{4\alpha},\qquad
  \beta_{,\tau,\tau}=\gamma_{,\tau,\tau}=\frac{1}{2}\mathrm{e}^{4\alpha}.
\end{equation}

这些方程的解应该描述度规从“初始”状态的演化\footnote{提醒一下，我们研究的是 $t\to0$ 时度规的演化；因此“初始”条件对应着更晚一些，而不是更早一些的时间.}，在初始状态中，方程的解通过选取一定的指数（并且 $p_{l}<0$），用公式（118.10）描述；设 $p_{l}=p_{1},\ p_{m}=p_{2},\ p_{n}=p_{3}$，于是
\begin{equation*}
  a=t^{p_{1}},\qquad b=t^{p_{2}},\qquad c=t^{p_{3}}
\end{equation*}
（在下面得到的结果中，不失一般性，可以设定这些表达式中的比例系数等于1）.在这种情况下 $abc=t,\ \tau=\ln t+\mathrm{const}$，因此方程（118.14）的初始条件可表示为下列形式：
\begin{equation*}
  \alpha_{,\tau}=p_{1},\qquad
  \beta_{,\tau}=p_{2},\qquad
  \gamma_{,\tau}=p_{3}.
\end{equation*}

方程组（118.14）的第一个方程具有粒子在指数形势垒的场中的一维运动方程的形式，其中 $\alpha$ 相当于坐标.在这个类比中，初始的卡兹涅尔状态对应于恒定速度为 $\alpha_{,\tau}=p_{1}$ 的自由运动.从势垒上反射后，粒子将再次以符号相反的速度 $\alpha_{,\tau}=-p_{1}$ 做自由运动.同时基于全部的三个方程（118.14）可发现
\begin{equation*}
  \alpha_{,\tau}+\beta_{,\tau}=\mathrm{const},\qquad
  \alpha_{,\tau}+\gamma_{,\tau}=\mathrm{const},
\end{equation*}
所以我们求得 $\beta_{,\tau}$ 和 $\gamma_{,\tau}$ 取数值 $\beta_{,\tau}=p_{2}+2p_{1},\ \gamma_{,\tau}=p_{3}+2p_{1}$.由此使用（118.6）确定 $\alpha,\beta,\gamma$，然后还有 $t$，我们得到
\begin{equation*}
  \mathrm{e}^{\alpha}\sim\mathrm{e}^{-p_{1}\tau},\qquad
  \mathrm{e}^{\beta}\sim\mathrm{e}^{(p_{2}+2p_{1})\tau},\qquad
  \mathrm{e}^{\gamma}\sim\mathrm{e}^{(p_{3}+2p_{1})\tau},\qquad
  t\sim\mathrm{e}^{(1+2p_{1})\tau},
\end{equation*}
即 $a\sim t^{p_{l}},\ b\sim t^{p_{m}},\ c\sim t^{p_{n}}$，其中
\begin{equation}
  p_{l}=\frac{|p_{1}|}{1-2|p_{1}|},\qquad
  p_{m}=-\frac{2|p_{1}|-p_{2}}{1-2|p_{1}|},\qquad
  p_{n}=\frac{p_{3}-2|p_{1}|}{1-2|p_{1}|}.
\end{equation}

这样一来，扰动的作用导致一个“卡兹涅尔状态”替换为另一个，并且 $t$ 的负幂指数从方向 $\boldsymbol{l}$ 跳到方向 $\boldsymbol{m}$：就是说如果过去 $p_{l}<0$，那么现在 $p_{m}<0$.在替换的过程中函数 $a(t)$ 经过最大值，而 $b(t)$ 经过最小值：先前减小的量 $b(t)$ 开始增大，先前增大的量 $a(t)$ 开始减小，而函数 $c(t)$ 则持续减小.先前增长的扰动本身（方程（118.7）中的 $a^{4}$ 项），开始下降和停息.度规的进一步演化以类似方式（由于（118.7）式中的 $b^{4}$ 项）导致扰动增长，变换到下一个卡兹涅尔状态，等等.

指数替换规则（118.15）适合于利用参数化的方式（118.12）表达：如果
\begin{equation*}
  p_{l}=p_{1}(u),\qquad p_{m}=p_{2}(u),\qquad p_{n}=p_{3}(u),
\end{equation*}
那么
\begin{equation}
  p_{l}=p_{2}(u-1),\qquad
  p_{m}=p_{1}(u-1),\qquad
  p_{n}=p_{3}(u-1).
\end{equation}

两个正指数当中较大的那一个仍然是正的.

在卡兹涅尔状态的更替过程中，存在着理解逼近奇点时度规演化特征的关键.

连续的更替（118.16）（这种更替伴随着负幂指数 $p_{1}$ 在 $\boldsymbol{l}$ 和 $\boldsymbol{m}$ 之间的跳跃）会一直持续到 $u$ 的初值的整数部分尚未消失，$u<1$ 的情况尚未出现的时候.数值 $u<1$ 按照（118.13）可以变换为 $u>1$；在这个时刻指数 $p_{l}$ 或者 $p_{m}$ 为负，而 $p_{n}$ 成为两个正数当中较小的一个（$p_{n}=p_{2}$）.下一个系列的更替将在方向 $\boldsymbol{n}$ 和 $\boldsymbol{l}$ 之间或者 $\boldsymbol{n}$ 和 $\boldsymbol{m}$ 之间交换负幂指数.在 $u$ 为任意（无理数的）初值时，更替过程会无限持续下去.

在对方程严格求解的时候，指数 $p_{1},p_{2},p_{3}$ 当然会失掉自己的严格含义.我们指出，这个状况会给这些数（以及同它们相关的参数 $u$）的确定带入某种“模糊性”，尽管这种模糊性很小，却会使得对某种方式挑选出的（例如，有理数的）数值 $u$ 的研究失去意义.正因为如此，只有在任意的、无理数的 $u$ 这种普遍情况下的规律性才具有现实意义.

这样一来，模型朝向奇点的演化过程由连续的振荡系列构成，在每一个振荡系列中沿着两个空间轴方向上的距离在振动，而沿着第三个轴的距离单调减小；体积按照近似于 $\sim t$ 的规律减小.在从一个系列变到下一个系列时，距离单调下降的方向会从一个轴更替到另一个轴.这些更替的顺序渐近地具有随机过程的特征.连续的振动系列的长度（就是说，每一个系列的“卡兹涅尔时代”数）\footnote{如果参数 $u$ 的“初始”值为 $u_{0}=k_{0}+x_{0}$（其中 $k_{0}$ 为整数，而 $x_{0}<1$），那么第一个振动系列的长度将为 $k_{0}$，而下一个系列的 $u$ 的“初始”值将为 $u_{1}=1/x_{0}\equiv k_{1}+x_{1}$，等等. 由此容易得出结论，连续系列的长度由单元 $k_{0},k_{1},k_{2},\cdots$ 给出，后者是将 $u_{0}$ 按照下式分解成无穷连分数（在 $u_{0}$ 为无理数的条件下）的单元
\[
u_{0}=k_{0}+\cfrac{1}{k_{1}+\cfrac{1}{k_{2}+\cfrac{1}{k_{3}+\cdots}}}.
\]
这个展式的更多连续单元值的分布符合统计学规律.}的交替顺序也会有随机性质.

连续的振动系列在逼近奇点的时候会变得稠密.在世界时间的任何一个有限时刻 $t$ 和时刻 $t=0$ 之间包含无穷多组振动.描述这个演化的时间进程的自然变量并不是时间 $t$ 本身，而是它的对数 $\ln t$，沿着 $\ln t$，逼近奇点的整个过程被拉伸至 $-\infty$.

在已表述的解中，我们从一开始就对问题做了些简化，即假定（116.3）中的矩阵 $\eta_{ab}(t)$ 是对角的.向度规中引入非对角的分量 $\eta_{ab}$ 不会改变已描述的度规演化的振动特征，以及连续的卡兹涅尔时代的指数 $p_{l},p_{m},p_{n}$ 的更替规律（118.16）.但是这会导致出现附加的性质：指数的更替伴随着它们所对应的坐标轴方向的改变\footnote{关于这一点，以及该类型均匀宇宙模型性质的其他细节可参考 В. А. Белинский, Е. М. Лифшиц, И. М. Халатников//УФН. 1970. Т. 102. С. 463; Adv. in Phys. 1970. V. 19. P. 525; ЖЭТФ. 1971. Т. 60. С. 1969.}.

\section{爱因斯坦方程一般宇宙学解的时间奇异性}

已经指出过，不能以弗里德曼模型对描述宇宙现今状态的适用性为依据，期望它同样适合于描述宇宙演化的早期阶段.与此相关首先会产生这样一个问题，存在时间奇点在多大程度上是宇宙学模型必有的一般性质，它和作为宇宙学模型基础的特殊简化假设（首先是对称性假设）到底有没有联系？我们强调，提到奇点时，我们是指物理奇点——物质密度和四维曲率张量的不变量趋于无穷大之处.

如果奇点的存在与这些假设无关，这就意味着，它不仅存在于爱因斯坦方程的特解中，而且也存在于通解中.对解的普适性的判据是看它包含的“物理上任意”的函数的数目.在通解中这些函数的数目应该足够用来在任意选择的时刻任意地给定初始条件（对于真空，这个数目是4；对于充满了物质的空间，该数目为8，参考§95\footnote{但是我们随即强调，对于爱因斯坦方程这样的非线性微分方程组，通解的概念不是单一的. 原则上可以存在多于一个的一般积分，其中每一个并不涵盖能想到的初始条件的所有的多样性，而只是它有限的一部分. 带有奇异性的通解的存在不排除同时存在另一些不具有奇异性的通解. 例如，没有理由怀疑存在这样的不带有奇异性的通解，它描述质量不太大的稳定的孤立物体.}）.

在整个空间和全部时间内寻求严格形式的通解显然是不可能的.但是对于解决提出的问题而言，只要研究奇异性附近的解的形式就足够了.

弗里德曼解具有的奇异性特征在于，空间距离的趋零按照同样的规律发生在所有的方向上.这种奇异性的类型不是足够一般的：它本质上是仅含坐标的三个任意函数的一类解（参考§113的习题）.同时我们指出，这些解只存在于填充了物质的空间.

在上一节中研究的振动类型的奇异性具有如下普适特征——存在带有这种奇异性的爱因斯坦方程的解，且它包含了所需的全部任意函数.这里我们简短地描述构造这种解的方法，而不深入计算的细节\footnote{它们可以在这些文章中找到：В. А. Белинский, Е. М. Лифшиц, И. М. Халатников//ЖЭТФ. 1972. Т. 62. С. 1606; Adv. in Phys. 1982. V. 31. P. 639.}.

如同在均匀模型中那样（§118），逼近奇点的状态由相互更替的“卡兹涅尔时代”的连续系列组成.在每个这样的时代的发展过程中，（在同步参考系中）空间度规张量的主项（对 $1/t$）具有（118.1）的形式，并且包含来自（118.10）的时间函数 $a,b,c$，但此时矢量 $\boldsymbol{l},\boldsymbol{m},\boldsymbol{n}$ 是空间坐标的任意函数（而不像在均匀模型中那样是完全确定的）.此时 $p_{l},p_{m},p_{n}$ 同样是函数（而不只是数），它们还像从前那样通过关系式（118.11）那样相互关联.以这种方式构造的度规在某个有限的时间段内满足真空中的场方程 $R_{0}^{0}=0$ 和 $R_{\alpha}^{\beta}=0$.方程 $R_{\alpha}^{0}=0$ 会导出三个关系式（不包含时间），它们应该添加到包含在 $\gamma_{\alpha\beta}$ 里的空间坐标的任意函数上.这些关系式将10个不同的函数相互联系起来：三个矢量 $\boldsymbol{l},\boldsymbol{m},\boldsymbol{n}$ 各自的三个分量和位于时间的指数上的一个函数（因为三个函数 $p_{l},p_{m},p_{n}$ 通过两个条件（118.11）相联系）.在确定物理上任意的函数的数目时应该同时考虑到，同步参考系还允许在不影响时间的前提下对三个空间坐标进行任意变换.因此度规总共包含 $10-3-3=4$ 个任意函数——这正是真空中场的通解应当包含的数目.

从一个卡兹涅尔状态到另一个的更替（就像在均匀模型中那样）源于六个方程 $R_{\alpha}^{\beta}=0$ 当中的三个里面存在这样一些项，它们在 $t$ 减小时增长得比其他项快.这样一来，它们起到破坏卡兹涅尔状态的扰动的作用.一般情况下，这些方程与方程（118.14）在形式上的区别仅在于位于它们右边的那些依赖于空间坐标的因子
$(\boldsymbol{l}\cdot\operatorname{rot}\boldsymbol{l}\,/\,\boldsymbol{l}\cdot\boldsymbol{m}\times\boldsymbol{n})^{2}$
（意味着，三个指数 $p_{l},p_{m},p_{n}$ 当中 $p_{l}$ 为负值）\footnote{对于均匀模型，这个因子与结构常数 $C^{11}$ 的平方相同并且按照定义是常数.}.但是由于方程（118.14）是相对于时间的常微分方程组，这个差别无论怎样也不会影响到它们的解以及由这个解得出的卡兹涅尔指数（118.16）的更替规律，还有在§118节中叙述的所有的更进一步的推论\footnote{如果对解中的任意函数加上一个附加条件 $\boldsymbol{l}\cdot\operatorname{rot}\boldsymbol{l}=0$，那么振荡会消失，卡兹涅耳状态将持续到点 $t=0$ 本身. 然而这样的解比一般情况下要求的少包含一个任意函数.}.

解的普适性程度在引入物质时不会降低：物质连同因它引入的全部4个新的坐标函数一起被“写入”度规，这4个新的坐标函数对于给出物质密度和三个速度分量的初始分布是必需的.物质的能量动量张量 $T_{i}^{k}$ 将 $1/t$ 的更高阶项（高于主项）带入场方程中（完全类似于§117节的习题3中针对平直均匀模型所示的那样）.

这样一来，时间奇点的存在是爱因斯坦方程解的相当普遍的性质，并且逼近奇点的状态在一般情况下具有振荡的特征\footnote{在爱因斯坦方程的通解中存在奇点的事实最先被彭罗斯（R. Penrose, 1965）采用拓扑学的方法所证明，然而，这些方法不能提供确定奇异性的具体解析特征的可能性. 这些方法的描述和借助于这些方法得到的定理可参考 R. Penrose. \textit{Structure of Space-Time}, W. A. Benjamin, N. Y., 1968.}.我们强调，这个特征与物质的存在无关（因而与物质的物态方程也无关），而是空的时空所固有的特性.弗里德曼解所固有的、与物质的存在有关的、单调各向同性类型的奇异性只具有特殊的意义.

在提到宇宙学意义上的奇异性时，我们指的是整个空间达到的奇点，而不是空间的有限部分（就像有限物体发生引力坍缩那样）达到的奇点.但是振动解的普适性给出了下面这个假设的根据，就是有限物体在共动参考系中坍缩到事件视界以内时达到的奇异性也具有同样的特征.

我们一直说逼近奇点的方向就是时间减小的方向；但是鉴于爱因斯坦方程相对于时间反演的对称性，那就可以同样合理地沿着时间增加的方向讨论逼近奇点的问题.然而，实际上，鉴于未来和过去在物理上并非等效，针对所提出问题的本身而言，在这两种情况之间就存在着实质性的差别.只有当未来的奇异性在之前某个时刻给定的任意初始条件下都可以实现的情况下，该奇异性才具有物理意义.可以明白，没有任何根据能够说明，在宇宙演化过程中的某个时刻达到的物质和场的分布能够对应于一些特殊的条件，而这些条件恰好是实现爱因斯坦方程的某个特解所要求的.

在以前所做的仅仅基于一些引力方程的研究中，关于奇异性的类型问题，通常未必能够给出唯一涵义的答案.自然地想到，和现实世界相对应的解的选择与某些深刻的物理条件有关，这些条件的建立仅仅基于一个现存的引力理论是不可能的，只有继续对物理理论进行整合才有可能揭示这些条件.就这个意义而言，原则上可以说，这种选择对应着某种特殊的（例如各向同性）类型的奇点.

最后，还有必要作出如下评论.爱因斯坦方程本身的适用范围绝不会在小的尺度或者大的物质密度这些方面受到限制，这是因为方程在这样的极限下不会导致任何内部矛盾（这与，例如，经典的电动力学方程不同）.在这个意义上，基于爱因斯坦方程研究时空度规的奇异性是完全合理的.然而，不用怀疑的是，实际上在该范围内应该存在重要的量子现象，关于后者，在目前的理论状况下，我们还没有什么可说.只有将来整合了引力理论和量子理论以后，才能够揭示经典理论的结论当中哪些仍然是有意义的.同时毋庸置疑的是，在爱因斯坦方程的解中，奇异性的产生（无论在宇宙学方面，还是有限物体的坍缩）这个事实本身有着深刻的物理涵义.不应忽视之点还有，在引力坍缩过程中那些极高密度的获得（在这种情况下还没有根据怀疑经典引力理论的合理性），足以用来讨论物理上“奇异”的现象.

\chapter*{译后记}

早在1959年，任朗、袁炳南二位先生就按1948年俄文第二版将本书译为中文出版，深受国内广大读者欢迎.半个多世纪以来，本书俄文原版历经数次修改和增补，和最初的版本相比篇幅已大为增加，全书的内容（特别是有关引力场理论的内容）多有更新，迫切需要出版包括这些修改和增补的新的中译本.这次的翻译是在任、袁二位先生原译本的基础上，按2006年俄文第八版进行的.

新译本的1—9章和12章由国家天文台邹振隆先生依据俄文第八版并参考英文第四版进行了校订.第10, 11, 13, 14章的增补部分、序言和索引由鲁欣依俄文第八版翻译，亦由邹振隆先生进行了校订.

特别要感谢北京师范大学物理系裴寿镛和赵峥二位先生，他们在出版前分别审读了1—9章和10—14章的译稿，提出了许多宝贵的意见，使得译稿更加完善.也要感谢北京大学力学系李植老师和中国科学院理论物理研究所刘寄星先生依据俄文版对序言译文所作的修改.

还值得一提的是，在本书的编辑过程中，部分读者通过本书的微博对部分译稿进行了试读，特别是复旦大学物理系施郁教授研究组博士生戴越同学通读了全稿，从科学性和文字两方面均提出了很多中肯的修改建议，使我获益匪浅.最后还要感谢高等教育出版社自然科学学术著作分社王超编辑耐心和细致的编辑.

受专业水平和语言水平的限制，译文中错误和不当之处在所难免.恳切地希望各位读者发现后批评指正，以便重印时修改.

\begin{flushright}
中国科学院物理研究所\quad 鲁欣\\
luxin@iphy.ac.cn\\
2012年6月于北京
\end{flushright}
